\documentclass[a4paper,10pt]{book}
%\documentclass[a4paper,10pt]{article}
\usepackage[utf8]{inputenc}
\usepackage[spanish]{babel}
\usepackage{showkeys}
\usepackage{bbm}
\usepackage[margin=2.5cm]{geometry}
\usepackage{amsmath,amssymb}
\usepackage{bigints}
\usepackage{enumerate} 
\usepackage{setspace}
\usepackage{hyperref}
\hypersetup{
    %bookmarks=true,         % show bookmarks bar?
    unicode=false,          % non-Latin characters in Acrobat’s bookmarks
    pdftoolbar=true,        % show Acrobat’s toolbar?
    pdfmenubar=true,        % show Acrobat’s menu?
    pdffitwindow=false,     % window fit to page when opened
    pdfstartview={FitH},    % fits the width of the page to the window
    pdftitle={Problema del TC},    % title
    pdfauthor={Ivan Mandelman},     % author
    pdfsubject={Tesis de Lic. en Matematica},   % subject of the document
    pdfcreator={Ivan Mandelman},   % creator of the document
    pdfproducer={Ivan Mandelman}, % producer of the document
    pdfkeywords={Tesis} {Optimizacion} {TC}, % list of keywords
    pdfnewwindow=true,      % links in new window
    colorlinks=true,       % false: boxed links; true: colored links
    linkcolor=red,          % color of internal links
    citecolor=green,        % color of links to bibliography
    filecolor=magenta,      % color of file links
    urlcolor=blue           % color of external links (cyan)
}

%\newcommand{\xcal}[1]{\text{\usefont{OMS}{cmsy}{m}{n}#1}}
%\DeclareMathAlphabet{\mathpzc}{OMS}{cmex}{m}{n}
\DeclareMathAlphabet{\mathpzc}{OT1}{pzc}{m}{n}
%\DeclareMathAlphabet{\mathcalligra}{T1}{calligra}{m}{n}
%\newcommand{\mycal}{\usefont{OT1}{pzc}{m}{n}\fontsize{10}{12}\selectfont} 

\newcommand{\N}{\mathbbm{N}}
\newcommand{\R}{\mathbbm{R}}
\newcommand{\Rn}{\R^n}
\newcommand{\arr}{\rightarrow}
\newcommand{\bx}{\bar{x}}
\newcommand{\blambda}{\bar{\lambda}}
\newcommand{\bmu}{\bar{\mu}}
\newcommand{\cN}{\mathcal{N}}
\newcommand{\cT}{\mathcal{T}}
\newcommand{\cI}{\mathpzc{I}}
\newcommand{\cE}{\mathpzc{E}}
\newcommand{\cA}{\mathpzc{A}}
\newcommand{\la}{\langle}
\newcommand{\ra}{\rangle}
\newcommand{\deb}{\rightharpoonup}

\newtheorem{teo}{Teorema}[section]
\newtheorem{lem}[teo]{Lema}
\newtheorem{cor}[teo]{Corolario}
\newtheorem{prop}[teo]{Proposici\'on}
\newtheorem{defi}[teo]{Definici\'on}
%\newtheorem{defi}{Definici\'on}[chapter]
%\newtheorem{ejem}[teo]{Ejemplo}
\newcounter{ejem}[section]
\newenvironment{ejem}[1][]{\refstepcounter{ejem}\par\medskip
   \noindent \textbf{Ejemplo~\theejem. #1} \rmfamily}{\medskip}
\newenvironment{dem}{\vspace*{-0.1cm}{\it Demostraci\'on:}}{\hfill
$\blacksquare$\\} 

\title{Trayectorias de tiempo óptimo \\para circuitos cerrados}
\author{Iván Mandelman}

\begin{document}

\maketitle

\section*{Resumen}
    Las trayectorias de tiempo óptimo en circuitos cerrados son usadas, por ejemplo, para determinar el camino que debe seguir un vehículo en una carrera o el recorrido de un riel en una línea de producción. Desde el punto de vista matemático, se plantea un problema de minimización de un cierto funcional en el espacio de las curvas cerradas y sujeto a condiciones propias del problema.
    
    El objetivo de este trabajo es el de analizar y resolver el problema de optimización de encontrar una trayectoria óptima para un circuito cerrado, de manera que minimice el tiempo que se demora en recorrerlo. En el problema a analizar plantearemos un modelo que contemple ciertas restricciones de permanencia en el circuito y fricción correspondientes.
    
    En este trabajo se analizará el problema a resolver para lograr modelarlo matemáticamente. Luego se darán las herramientas teóricas y técnicas numéricas para resolver el problema. Con el modelo matemático listo, se procederá a realizar una discretización del problema para resolverlo computacionalmente y finalmente se exhibirán los resultados numéricos obtenidos con su interpretación gráfica correspondiente.

\tableofcontents

\pagebreak

\section*{Notaci\'on}
{\setlength{\parindent}{0cm}
%{\setstretch{1.5}
$:=$ igual por definición.

$\emptyset$ conjunto vacío.

\hangindent=1cm \textbf{p.p.} para casi todo punto, es decir, en todos los puntos excepto posiblemente en los de medida (medida dada por el contexto) cero.

$x\to f(x)$ mapea el punto $x$ hacia $f(x)$.

$F:X\arr Y$ la función $F$ sale del espacio $X$ y llega al espacio $Y$.

$\R$ los números reales.

$\Rn$ es el espacio euclídeo $n$-dimensional.

$\Rn_+=\{x\in\Rn\,|\,x_i\geq0,\,i=1,\ldots,n\}$ es el ortante no negativo.

\hangindent=1cm $L_2[0,1]$ el espacio de Hilbert de clases de equivalencias de valores reales, cuadrado integrable, funciones $\psi(t)$, con $\psi_1\sim\psi_2$ si $\psi_1(t)=\psi_2(t)$ para todo $t\in[0,1]$ excepto posiblemente en un conjunto de medida Lebesgue cero, y $\la\psi,\varphi\ra=\int_0^1\psi(t)\varphi(t)dt$.

\hangindent=1cm $L_p(\Omega)$ espacio de Banach de funciones medibles Lebesgue $\psi\!:\!\Omega\arr\R$ con norma finita $\|\psi\|_p\!:=\!(\int_{\Omega}\!|\psi(\omega)|^p d\omega)^{\frac{1}{p}}$.

$L_p(\Omega)_+\subset L_p(\Omega)$ el conjunto de funciones con valores no negativos p.p. $x\in\Omega$.

\hangindent=1cm $W^{m,s}(\Omega)=\{\psi\in L_s(\Omega)\,|\, D^q\psi\in L_s(\Omega)\mbox{ si }|q|\leq m\}$ espacio de Sobolev, donde $D^q\psi=\frac{\partial^{|q|}\psi}{\partial_{x_1}^{q_1}\ldots\partial_{x_l}^{q_l}}$ y $|q|=q_1+\ldots+q_l$.

$H^m(\Omega)$ es el espacio de Sobolev $W^{m,2}(\Omega)$.

$C_c^\infty(\Omega)$ son las funciones infinitamente diferenciables de soporte compacto.

\hangindent=1cm $\mathcal{L}(X,Y)$ espacio de operadores lineales continuos $A:X\arr Y$ con la norma $\|A\|=\sup_{\|x\|<1}\|Ax\|$ siendo $X$ e $Y$ espacios normados.

$X^*=\mathcal{L}(X,\R)$ espacio dual de $X$.

$\mbox{span}\{S\}$ es el subespacio generado por el conjunto $S\subset X$ donde $X$ es un espacio vectorial.

int($S)=\{x\in S\,|\, \exists\mbox{ un entorno }V \mbox{ de }x\mbox{ tal que }V\subset S\}$ es  el interior del conjunto $S$.

$r(t)=o(t)$ si $r:\R_+\to\R$ cumple que $r(t)/t\to 0$ cuando $t\to0$.

$\varliminf_{n\to\infty}x_n=\sup_{n\geq0}\mbox{inf}_{k\geq n}(x_k)$ es el límite inferior de una sucesión $\{x_n\}$.

$\varlimsup_{n\to\infty}x_n=\mbox{inf}_{n\geq 0}\sup_{k\geq n}(x_k)$ es el límite superior de una sucesión $\{x_n\}$.

$\mbox{cl}(S)$ es la clausura (con la topología usual) de el conjunto $S$.

Si $A$ y $B$ son espacios vectoriales, $A+B=\{a+b\mid a\in A,\, b\in B\}$ es la suma de Minkowski.

$A^t$ es la transpuesta de la matriz $A$.

$\nabla f(x)=(\partial f(x)/\partial x_1,\ldots,\partial f(x)/\partial x_n)$ gradiente de la función $f:\Rn\arr\R$ en el punto $x\in\Rn$.

$A(:,m:l)$ son todas las filas, y las columnas desde la $m$ hasta la $l$ de la matriz $A$.

$[A;B]$ es una matriz formada por el bloque $A$ y en la fila debajo, el bloque $B$.

$B(a,\varepsilon)=\{x\in\Rn\mid \|x-a\|_{\Rn}<\varepsilon\}$ es la bola abierta de $\Rn$.

$B[a,\varepsilon]=\{x\in\Rn\mid \|x-a\|_{\Rn}\leq\varepsilon\}$ es la bola cerrada de $\Rn$.

$\mathrm{Im}(f)$ es la imagen de la función $f$. 

En el caso en que el contexto no lo especifique, la derivada de una función en un punto $x$, $f'(x)$ con $f:\Rn\arr\R$, será considerada como una matriz de $n$ filas y $1$ columna. De la misma manera, consideraremos un vector en $\Rn$ como un vector columna.

Si no se aclara lo contrario $\{x_k\},\, k\in\N$ denotara la sucesión de números reales $x_1,x_2,\ldots$.

$C^l(\Omega)$ es el espacio de Banach de funciones $\psi:\Omega\arr\R,\ \Omega\in\Rn$. $l$-veces continuamente diferenciables. 

Si $A:X\arr Y$ es un operador lineal, $\mathrm{Ker}(A)=\{x\in X\mid Ax=0\}$ es el núcleo de $A$.

Si $b,d\in\Rn$, $b\leq d$ es $b_i\leq d_i,\,\forall i=1,\ldots,n$.

%poner toda la notacion
}

\chapter*{Introducción}

El objetivo de este trabajo resolver el problema de optimizacón de encontrar la trayectoria óptima para un circuito cerrado, de manera que se minimice el tiempo que se tarda en recorrerlo.

En el Capítulo 1, describimos el problema que nos interesa resolver y realizamos una formulación matemática del mismo, así, obtenemos un problema de optimización con restricciones y es posible utilizar marco teórico matemático para resolverlo. Éste marco teórico se da a continuación en el Capítulo 2. De esta forma, terminamos de formular el problema de optimización en dimensión infinita en el Capítulo 3, cuya idea es obtener un problema como el siguiente

\begin{equation*}
\begin{array}{rl}
\displaystyle \mathop{\mathrm{minimizar}}_{x} & \displaystyle f(x) \\
\mathrm{sujeto \ a} & G(x)\in K,\, x\in C
\end{array}
\end{equation*}
donde detallamos cuáles son los espacios y funciones constituyentes del problema.




Luego, en el Capítulo 4, procedemos a formular el problema en dimension finita (e incluímos el marco teórico necesario), de manera que podamos utilizar algún método computacional para obtener resultados numéricos, lo que podemos ver en el Capítulo 5 y el Capítulo 6. En este último, podremos observar los resultados en gráficas descriptivas que muestran las trayectorias solución obtenidas.

%explicar el motivo de la tesis con palabras (sin formulas)


\chapter{La trayectoria óptima}

\section{El ``Problema del TC''}
Supongamos que nos encontramos en alguno de los circuitos del TC (turismo carretera), estos circuitos (o pistas), tienen un punto (o línea) de salida, que es el mismo que el de llegada, y el objetivo, es dar una determinada cantidad de ``vueltas'' en el menor tiempo posible. Es sensato entonces, decir que, encontrando una estrategia óptima para \emph{una} vuelta, estaremos cumpliendo el objetivo antes mencionado.

Para no ser restrictivos con que el vehículo para dar la vuelta es un auto, pensaremos que nos movemos ``puntualmente''.

Además, hay otras consideraciones a tener en cuenta:
\begin{itemize}
 \item debemos mantenernos dentro de la pista, esto es indispensable pues se plantea que el terreno fuera de la pista es infranqueable en muchos casos.

 \item La fricción de la pista, que regula los límites de la aceleración, y también, la trayectoria.
\end{itemize}
Por lo tanto la idea es, minimizar el tiempo en que damos la vuelta, teniendo en cuenta las consideraciones previas.
Para resolver esto, planteamos un problema de minimización, con determinadas restricciones. La función a minimizar es el tiempo, y algunas de las restricciones deberán ser la permanencia en pista y la fricción.
%\section{formulacion matematica}
\section{Formulacion matemática}
%\chapter{preliminares de optimizacion}
%\section{dim infinita}
%% finalizar planteando el problema en este contexto definiendo las funciones y sus derivadas
%\section{discretizacion y dim finita}
%% primero en abstracto y luego
%%discretizacion de los espacios del problema y de las funciones (y derivadas)
%\chapter{resolucion numerica}
%\section{metodo computacional}
%\section{resultado numerico}

El objetivo ahora es definir la función \textbf{tiempo final} $t_f$ que queremos minimizar, sujeta a las dos condiciones mencionadas anteriormente. Matem\'aticamente la condici\'on de mantenerse dentro de la pista puede escribirse como
\[
x(t)\in \mathcal{P}, \quad\forall t\in[0,t_f],
\]
donde $x$ es la trayectoria del vehículo y $\mathcal{P}$ es el conjunto de curvas que se mantienen dentro de la pista. La condici\'on de fricción regula las aceleraciones tangencial $a_t$ y normal $a_n$ para que cumplan
\begin{equation}\label{elipse}
 \left(\frac{a_t}{\mu_t}\right)^2+ \left(\frac{a_n}{\mu_n}\right)^2\leq g^2, 
\end{equation}
donde $g$ es la acelaración de la gravedad, $\mu_t$ y $\mu_n$ son los coeficientes de rozamiento tangencial y normal, respectivamente. Tal condici\'on es conocida en la litaratura como \textbf{elipse de fricción}.

Para estudiar este problema, definiremos la curva $x$ a través de sus funciones de movimiento. Una herramienta que facilitará las cuentas es reparametrizar la curva por longitud de arco.
Entonces consideramos la función \textbf{trayectoria} de la curva como $x:[0,t_f]\mapsto \R^2$ tal que
\begin{equation}\label{trayectoria}
x(t) = p c(S(t)), 
\end{equation}
donde \(t_f\) es el tiempo final, $c:[0,1]\arr\R^2$ es una curva parametrizada por longitud de arco, es decir, para $s\mapsto c(s)$ con $s\in[0,1]$
\begin{equation}\label{pla}
\left \|\frac{dc}{ds}(s)\right\|_{\R^2}=1,    
\end{equation}
la funci\'on $S:[0,t_f]\arr [0,1]$ es la longitud de arco normalizada y $p\in(0,+\infty)$ es la longitud de $x$, pues
\[
 \int_0^{t_f} \left\|\frac{dx}{dt}(t)\right\|_{\R^2} dt = p \int_0^{t_f} \frac{dS}{dt}(t)dt=p(S(t_f)-S(0))=p.
\]


Para definir la funci\'on $t_f$ usaremos las propiedades de las funciones anteriormente dadas. Para ello definimos la función \textbf{velocidad} $\nu:[0,t_f]\arr\R^2$ tal que
\begin{equation}
  \nu(t)=\frac{dx}{dt}(t)=p\frac{dc}{ds}(S(t))\frac{dS}{dt}(t), \label{velocidad}
\end{equation}
y la funci\'on rapidez (magnitud de la velocidad) $\tilde{v}:[0,t_f]\arr\R$ tal que
\[
 \tilde{v}(t)=p\frac{dS}{dt}(t).
\]
Note que estas funciones estan definidas en un intervalo de tiempo que depende de $t_f$.

De la definici\'on de longitud de arco, se puede demostrar que $S$ tiene inversa $T:[0,1]\arr [0,t_f]$ con $T(0)=0$ (una explicación detallada puede verse en el Capítulo \ref{reparametr}), usando esto
definimos una nueva funci\'on de rapidez $v:[0,1]\arr\R$ tal que
\begin{equation}\label{v}
v(s)=\tilde v(T(s))=p\frac{dS}{dt}(T(s)). 
\end{equation}
Como $s=S(T(s))$, derivando en relaci\'on a $s$ obtenemos 
\begin{equation}\label{auxdSdt}
1=\frac{dS}{dt}(T(s)) \frac{dT}{ds}(s)= \frac{1}{p} v(s)\frac{dT}{ds}(s), 
\end{equation}
por lo tanto
\[
 \frac{dT}{ds}(s)=\frac{p}{v(s)}.
\]
Integrando, se obtiene que el tiempo final es
\begin{equation}\label{tiempo}
t_f=T(1)-T(0)=\int_{0}^{1} \frac{dT}{ds}(s) ds= \int_{0}^{1} \frac{p}{v(s)}ds.    
\end{equation}


Para escribir la condici\'on de la elipse de fricción en t\'ermino de estas mismas funciones, calcularemos la acelaración de la trayectoria
\begin{equation}\label{aceleracion}
\frac{d^2 x}{dt^2}(t)= p\frac{d^2 c}{ds^2}(S(t)) \left(\frac{dS}{dt}(t)\right)^2 + p\frac{dc}{ds}(S(t))\frac{d^2 S}{dt^2}(t),
\end{equation}
donde en el primer t\'ermino se encuentra la informaci\'on de la aceleraci\'on normal y en segundo el de la aceleraci\'on tangencial. Esto puede demostrarse usando que $c$ est\'a parametrizada por longitud de arco y por lo tanto
\[
 0=\frac{d}{ds}1=\frac{d}{ds}\left\| \frac{dc}{ds}(s)\right\|_{\R^2}^2=2 \left\langle\frac{dc}{ds}(s) , \frac{d^2 c}{ds^2}(s)\right\rangle. 
\]
Por lo tanto, la aceleración normal tiene magnitud
\begin{equation}\label{acelnormal}
a_n(s)=p\left(\frac{dS}{dt}(T(s))\right)^2\left\| \frac{d^2 c}{ds^2}(s) \right\|=\frac{1}{p} v(s)^2\left\| \frac{d^2 c}{ds^2}(s) \right\|_{\R^2},
\end{equation}
donde usamos \eqref{v}. Por otro lado, usando \eqref{v} y (\ref{auxdSdt}),
\[
 \frac{dv}{ds}(s)=p\frac{d^2 S}{dt^2}(T(s))\frac{dT}{ds}(s)=p\frac{d^2 S}{dt^2}(T(s))\frac{p}{v(s)}.
\]
Por lo tanto la aceleración tangencial es
\[
 a(s):=a_t(s)=p\frac{d^2 S}{dt^2}(T(s))=\frac{1}{p} v(s)\frac{d v}{ds}(s)=\frac{1}{p}\frac{d}{ds}\left(\frac{1}{2}v(s)^2\right)
\]
Con esto, obtenemos que
\begin{equation}\label{rel}
p \int_0^s a(r)dr =\frac{1}{2}(v(s)^2-v(0)^2). 
\end{equation}


Para describir el conjunto $\mathcal{P}$ (curvas que se mantienen dentro de la pista) en t\'erminos de las funciones anteriores, considere que $\alpha:[0,1]\arr\R^2$ y $\beta:[0,1]\arr\R^2$ son las curvas que describen la calzada interna y externa de la pista, respectivamente, y sea $\kappa:[0,1]\arr(0,+\infty)$ el ancho de la calzada, o sea $\kappa(s)=\|\beta(s)-\alpha(s)\|_{\R^2}$. Pediremos que para todo $s\in[0,1]$
\begin{equation}\label{pista2}
\begin{array}{c}
 \|x(T(s))-\alpha(s)\|_{\R^2}^2 \leq \kappa(s)^2,\\[2mm]
 \|x(T(s))-\beta(s)\|_{\R^2}^2 \leq \kappa(s)^2.
\end{array}
\end{equation}
Observemos que lo que describe este sistema de desigualdades es lo siguiente: la distancia euclídea entre la trayectoria y la calzada interna, y a la vez entre la trayectoria y la calzada externa deben ser menores o iguales que el ancho de la pista. Entonces, si fuese que la trayectoria se desvía fuera de la pista hacia el interior del circuito, \'esta no cumpliría la restricción que considera la calzada externa (pues $\|x(T(s))-\beta(s)\|_{\R^2}^2 > \kappa(s)^2$), y análogamente si la trayectoria se desvía fuera de la pista hacia el exterior del circuito ($\|x(T(s))-\alpha(s)\|_{\R^2}^2 > \kappa(s)^2$).
Por lo tanto, con estas restricciones nos aseguramos la permanencia en pista.

También, cabe destacar que al origen de coordenadas lo ubicaremos siempre dentro del área encerrada por la curva interna $\alpha$, lo cual podemos describir mediante las siguientes ecuaciones.
\begin{equation*}
    \begin{array}{c}
         \|\beta(s)\|_{\R^2}>\kappa(s),\,\forall s\in[0,1],\\
         \|\beta(s)\|_{\R^2}>\|\alpha(s)\|_{\R^2},\,\forall s\in[0,1].
    \end{array}
\end{equation*}

Por último, como nuestras curvas (interna y externa) intentarán describir un circuito plano y cerrado, es necesario que no se intersequen consigo misma, esto es, las curvas serán simples (ver Apéndice, definición \eqref{curvasimple}).
%\begin{equation}\label{pista1}
%\begin{array}{c}
% (x(T(s))-\alpha(s))\cdot n_\alpha(s) \leq 0,\\
% (x(T(s))-\beta(s))\cdot n_\beta(s) \leq 0.
%\end{array}
%\end{equation}




% Para limitar las vueltas que puede dar la curva $c$ usaremos su número de vueltas $\gamma(c)$, i.e.,
% \begin{eqnarray*}
%  \gamma(c)&=&\frac{1}{2\pi i}\int_c \frac{dz}{z}=\frac{1}{2\pi i}\int_c \frac{x dx + y dy + i(-y dx + x dy)}{x^2+y^2}\\
%  &=&\frac{1}{2\pi}\int_c \frac{-y dx + x dy}{x^2+y^2}=\frac{1}{2\pi}\int_0^1 \frac{-c_2(t) c_1'(t) + c_1(t) c_2'(t)}{c_1(t)^2+c_2(t)^2}dt,
% \end{eqnarray*}
% donde se usó que $c$ es una curva cerrada y por lo tanto
% \begin{eqnarray*}
%  \int_c \frac{x dx + y dy }{x^2+y^2}&=&\int_0^1 \frac{c_1(t)c_1'(t)+c_2(t)c_2'(t)}{\|c(t)\|^2}dt=\frac{1}{2}\int_0^1\log(\|c(t)\|^2)dt\\
%  &=&\frac{1}{2}(\log(\|c(1)\|^2)-\log(\|c(0)\|^2))=0.
% \end{eqnarray*}

 

%\textbf{Hablar de las otras restricciones!!!}

\subsection{Problema de optimización}
Estamos listos ahora para comenzar a escribir el problema de optimización.
Recapitulando lo hecho anteriormente, identificamos las incógnitas:
\begin{itemize}
 \item la curva $c$,
 \item la rapidez $v$,
 \item la aceleración tangencial $a$,
 \item la longitud de la curva $p$, y
 \item la velocidad inicial $v(0)$ a la cual llamaremos $v_0$.
\end{itemize}

Por \eqref{tiempo} ya tenemos definida la función objetivo a minimizar.

De las cuentas anteriormente realizadas surjen las restricciones que necesitamos:
\begin{enumerate}
\item la parametrización por longitud de arco de $c$ (\ref{pla}).
\item La curva $c$ debe ser cerrada, luego $c(1)=c(0)$.
\item La velocidad y aceleraci\'on de $c$ deben ser curvas cerradas, entonces $\frac{dc}{ds}(1)=\frac{dc}{ds}(0)$, $\frac{d^2c}{ds^2}(1)=\frac{d^2c}{ds^2}(0)$.

\item La condición de la elipse de fricción (\ref{elipse}), usando (\ref{acelnormal}).

\item La rapidez inicial y final deben ser iguales, por lo tanto $v(1)=v(0)$.

\item La aceleraci\'on tangencial inicial y final deben ser iguales, entonces $a(1)=a(0)$. 

\item La rapidez y aceleraci\'on tangencial deben satisfacer la condición (\ref{rel}).

\item La permanencia en pista (\ref{pista2}), teniendo en cuenta (\ref{trayectoria}).
\end{enumerate}

Para formular el problema de acuerdo al est\'andar en optimizaci\'on continua, escribiremos las restricciones de igualdad con un igual a cero y las de desigualdad con un menor o igual a cero. Por lo tanto el planteo del problema queda de la siguiente manera:
\[
\begin{array}{rl}
\displaystyle \mathop{\mathrm{minimizar}}_{(c,a,v,p,v_0)} & \displaystyle \int_{0}^{1} \frac{p}{v(s)}ds \\
 \mathrm{sujeto \ a} & \left\|\frac{dc}{ds}(s)\right\|_{\R^2}^2-1=0,\,\forall s\in[0,1]\\[2mm]
                     & c(1)-c(0)=0,\\[2mm]
                     & \frac{dc}{ds}(1)-\frac{dc}{ds}(0)=0,\\[2mm]
                     & \frac{d^2c}{ds^2}(1)-\frac{d^2c}{ds^2}(0)=0,\\[2mm]
                     & v(1)-v_0=0,\\[2mm]
                     & a(1)-a(0)=0,\\[2mm]
                     &  \frac{1}{2}(v(s)^2-v_0^2)-p\int_0^s a(r)dr=0,\, \forall s\in[0,1],\\ [2mm]
                     & \frac{1}{\mu_t^2} a(s)^2 + \frac{1}{\mu_n^2 p^2}v(s)^4\left\| \frac{d^2 c}{ds^2}(s) \right\|_{\R^2}^2-g^2\leq 0,\,\forall s\in[0,1],\\[2mm]
                     &  \|pc(s)-\alpha(s)\|_{\R^2}^2-\kappa(s)^2 \leq 0,\,\forall s\in[0,1], \\[2mm]
                     &  \|pc(s)-\beta(s)\|_{\R^2}^2-\kappa(s)^2 \leq 0,\, \forall s\in[0,1].
  \end{array}
\]
donde $(c,a,v,p,v_0)\in H^2[0,1]^2\times L^2[0,1]\times L^2[0,1]\times\R_+\times \R_+$.

Note que la elección de los espacios tiene sentido, pues, la velocidad y aceleración deben ser integrables en $[0,1]$, y la curva (que tiene 2 coordenadas espaciales), por regularidad, debe poder ser derivable dos veces, de manera que el espacio $H^2$ nos provee esto, además de ser un espacio de Hilbert, lo cual será de mucha utilidad en el marco teórico.

%\textbf{arreglar esto}

Luego de obtener $\bar v$, $\bar c$, $\bar p$ y $\bar v_0$, solución del problema de optimización anterior, se define el tiempo final óptimo como
\[
 t_f=\int_{0}^{1} \frac{\bar p}{\bar v(s)}ds.
\]
Además, obtenemos la trayectoria óptima
\[
 x(T(s))=\bar p\bar c(s),\qquad s\in[0,1].
\]
% y se obtiene $S$ como solución de la ecuación diferencial
% \[
%  S'(t)=\frac{1}{\bar p}\bar v(S(t)), \quad S(0)=0,\ S(t_f)=1.
% \]
% Con esto se puede obtener la trayectoria óptima $x(t)=\bar p \bar c(S(t))$.

% Otro planteo del problema
% 
% Dada una función de curvatura $k:[0,1]\mapsto\R$ podemos definir
% \[
%  c(s)=\left(\int_0^s \cos(\theta(t))dt + x_0, \int_0^s \mathrm{sen}(\theta(t))dt + y_0\right),
% \]
% donde
% \[
%  \theta(t)=\int_0^t k(r) dr + \varphi,
% \]
% con $x_0,y_0,\varphi\in\R$. Note que con esta elección obtenemos que
% \[
%  \left\|\frac{dc}{ds}(s)\right\| =1, \qquad \left\|\frac{d^2c}{ds^2}(s)\right\|=|k(s)|.
% \]
% 
% 
% 
% \subsection{Otro problema de optimización}
% 
% \begin{eqnarray*}
%  \mathop{\mathrm{minimizar}}_{(a,v,k,x_0,y_0,\varphi)} && \displaystyle \int_{0}^{1} \frac{ds}{v(s)} \\
%  \mathrm{sujeto \ a} &&\textstyle \frac{1}{\mu_t^2} a(s)^2 + \frac{1}{\mu_n^2}v(s)^2k(s)^2-g^2\leq 0,\\
%                      &&\textstyle \int_0^1 \cos(\theta(t))dt=0,\\
%                      &&\textstyle \int_0^1 \mathrm{sen}(\theta(t))dt=0,\\
%                      && v(1)-v(0)=0,\\
%                      &&\textstyle  \frac{1}{2}(v(s)^2-v(0)^2)-\int_0^s a(r)dr=0,\\ 
%                      && h(c(s))\leq 0,\\
%                      && \gamma(c)-1=0,
% \end{eqnarray*}
% donde $\theta(t)=\int_0^t k(r) dr + \varphi$ y la pista se representa por $C=\{x\in\R^2\mid h(x)\leq 0\}$.


Una forma de simplificar el problema, es reducir la cantidad de variables, escribiendo a $v$ en función de $a$. Entonces, usando (\ref{rel}) 
\begin{equation}\label{vel}
v(s)=\sqrt{2p \int_0^s a(r)dr + v_0^2}
\end{equation}
Y el problema se reescribe como
\begin{equation}\label{problema}
\begin{array}{rl}
\displaystyle \mathop{\mathrm{minimizar}}_{(c,a,p,v_0)} & \displaystyle \int_{0}^{1} \frac{p}{\sqrt{2p \int_0^s a(r)dr + v_0^2}}ds \\
 \mathrm{sujeto \ a} & \left\|\frac{dc}{ds}(s)\right\|_{\R^2}^2-1=0,\\[2mm]
                     & c(1)-c(0)=0,\\[2mm]
                     & \frac{dc}{ds}(1)-\frac{dc}{ds}(0)=0,\\[2mm]
                     & \frac{d^2c}{ds^2}(1)-\frac{d^2c}{ds^2}(0)=0,\\[2mm]
                     & \int_{0}^{1} a(r)dr=0,\\[2mm]
                     & a(1)-a(0)=0,\\[2mm]
                     & \frac{1}{\mu_t^2} a(s)^2 + \frac{1}{\mu_n^2 p^2}(2p \int_0^s a(r)dr + v_0^2)^2\left\| \frac{d^2 c}{ds^2}(s) \right\|_{\R^2}^2-g^2\leq 0,\,\forall s\in[0,1],\\[2mm]
                     &  \|pc(s)-\alpha(s)\|_{\R^2}^2-\kappa(s)^2 \leq 0,\,\forall s\in[0,1], \\[2mm]
                     &  \|pc(s)-\beta(s)\|_{\R^2}^2-\kappa(s)^2 \leq 0, \,\forall s\in[0,1].

  \end{array}
\end{equation}
donde $(c,a,p,v_0)\in H^2[0,1]^2 \times L^2[0,1]\times \R_+\times \R_+$. 

En el siguiente capítulo, mostraremos las herramientas teóricas y técnicas numéricas usadas para resolver el problema \eqref{problema}.

\chapter{Preliminares de geometría diferencial}\label{reparametr}

En esta capítulo, veremos el marco teórico necesario para comprender la reparametrización por longitud de arco de una curva. Como sólo utilizamos curvas planas, adecuaremos la teoría a ello, pero es importante remarcar que se puede extender a curvas en el espacio, es decir, aquellas que llegan a $\R^3$. Para comenzar damos la definición de curva.
\begin{defi}
 Una curva (parametrizada diferenciable) es una función infinitamente diferenciable $c:I\arr\R^2$ donde $I\subset\R$ es un intervalo abierto.
\end{defi}

Como en este trabajo las curvas que consideramos son siempre suaves, se da la siguiente definición que evita que las curvas tengan ``picos''.

\begin{defi}
 Una curva $c:I\arr\R^2$ se dice regular si $c'(t)\neq0,\,\forall t\in I$. Los puntos $t\in I$ donde $c'(t)=0$ se denominan puntos singulares.
\end{defi}

\textbf{Observación}. Los puntos singulares son justamente aquellos donde podría no estar definida la recta tangente a la curva.

Ahora, definimos la siguiente función, que nos permite medir en cierto sentido la longitud de una curva.

\begin{defi}
 La longitud de arco de una curva $c:I\arr\R^2$ desde $t_0\in I$ es $$ s(t):=\int_{t_0}^t\|c'(t)\|_{\R^2}dt.$$
\end{defi}
De esta manera, $s(t)$ mide la curva entre $c(t_0)$ y $c(t)$. Además, se puede observar que $s(t_0)=0$.

Recorrer una curva con velocidad constante y normalizada, resulta muy útil para trabajos que necesitan analizar curvas, por lo que la siguiente definición es imprescindible.

\begin{defi}
 Decimos que una curva $c:I\arr\R^2$ está parametrizada por longitud de arco si $$\|c'(t)\|_{\R^2}=1,\,\forall t\in I.$$
\end{defi}

\textbf{Observación}. En consecuancia, $s(t)=t-t_0$, es decir, recorre $t$ unidades de medida en $t$ unidades de tiempo.

Finalmente, podemos observar que, dada una curva regular $c:I\arr\R^2$, siempre se la puede reparametrizar por longitud de arco, esto es, tomamos $\beta(s):=c(t(s)$, donde $t:s(I)\arr I$ es la inversa de la función longitud de arco $s$ (con un abuso de notación), la cual existe porque $s'(t)=\|c'(t)\|_{\R^2}>0$ y por lo tanto $s$ es creciente. En efecto $$\beta'(s)=c'(t(s))t'(s),$$
entonces $$\|\beta(s)\|_{\R^2}=\frac{1}{\|c'(t)\|_{\R^2}}\|c'(t)\|_{\R^2}=1,\,\forall t\in I.$$

\textbf{Observación}. Como dijimos antes, $s(t_0)=0$. Luego si $t$ es la función inversa de la longitud de arco $s$, obtenemos que $t(0)=t_0$.

Se puede ver con detalle la utilidad de la reprametrización por longitud de arco en \cite[Section 1-5.]{do1990geometria}.

A continuación veremos unos resultados fundamentales sobre curvas cerradas. Entonces definimos curva cerrada y curva simple.

\begin{defi}
 Sea $[a,b]\subset\R$, $c:[a,b]\arr\R^2$ una curva. Entonces $c$ se dice cerrada, si $c(a)=c(b)$. 
\end{defi}

\begin{defi}\label{curvasimple}
Sea $c:I\subset\R\arr\R^2$ una curva, entonces $c$ se dice simple si es inyectiva.
\end{defi}

Además necesitaremos definir la curvatura de una curva.

\begin{defi}
Sea $c:I\subset\R\arr\R^2$ una curva parametrizada por longitud de arco. El número $k(s):=\|c''(s)\|_{\R^2}$ se denomina curvatura de $c$ en $s$ (la curvatura de una curva plana lleva signo dependiendo la orientación, horaria o antihoraria, de la misma).
\end{defi}

Enunciamos entonces, el siguiente teorema, cuya demostración se puede observar en \cite[Theorem 1, pg 33-35]{do1990geometria}

\begin{teo}[Desigualdad isoperimétrica]\label{desigualdadisoperimetrica}
Sea $c$ una curva plana, simple y cerrada con longitud $p$, y sea $A$ el área de la región limitada por $c$. Entonces
\begin{equation*}
    p^2-4\pi A\geq0,
\end{equation*}
y la igualdad se obtiene si, y sólo si $c$ es una circunferencia.
\end{teo}

Además será de utilidad la siguiente proposición.

\begin{prop}\label{cotainfcurvatura}
 Sea $c:(a,b)\arr\R^2$ una curva parametrizada por longitud de arco. Asuma que existe un $t_0\in(a,b)$ tal que la distancia de $c$ al origen ($t\mapsto\|c(t)\|_{\R^2}$) alcanza su máximo en $t_0$. Entonces vale que 
 $$\|c''(t_0)\|_{\R^2}\geq\frac{1}{\|c(t_0)\|_{\R^2}}.$$
\end{prop}

\begin{dem}
Como $c$ está parametrizada por longitud de arco, entonces vale que $$\|c'(t)\|_{\R^2}=1\ \forall t\in(a,b),$$
ahora, si tenemos en cuenta a la función $f(t):=\|c(t)\|^2_{\R^2}$, sabemos que $f$ tiene un máximo en $t_0$ y esto implica que $$f''(t_0)=2(\la c(t_0),c''(t_0)\ra_{\R^2}+\|c'(t_0)\|_{\R^2})\leq 0,$$
y luego $$1\leq-\la c(t_0),c''(t_0)\ra_{\R^2}\leq\|c(t_0)\|_{\R^2}\|c''(t_0)\|_{\R^2}.$$
Finalizando la demostración.
\end{dem}

%\begin{prop}\label{desigualdadintegralcurvatura}
% Sea $c:[0,p]\arr\R^2$ una curva plana y cerrada, donde $p$ es la longitud de la curva. %Si $k$ es la curvatura de $c$, vale que 
% \begin{equation*}
%     \int_0^pk(s)ds=2\pi T
% \end{equation*}
% donde $T$ es un entero (que valdrá $1$ o $-1$ dependiento el sentido de la curva) %denominado índice de rotación.
%\end{prop}

%El resultado anterior puede encontrarse en \cite[Pg 36,37,38]{do1990geometria}.


\chapter{Preliminares de optimización}
\section{Dimensión infinita}
\subsection{Existencia de soluciones}\label{existenciadesoluciones}

Para saber si un problema de optimización, como el \eqref{problema} planteado en el capítulo anterior, tiene solución necesitamos desarrollar varios conceptos teóricos. Al trabajar con funcionales lineales continuos, denotaremos por $\la\cdot,\cdot\ra_{X^*,X}$ a la acción dual, i.e., $\la f,x\ra_{X^*,X}=f(x)$ $\forall\,f\in X^*$ $\forall\, x\in X$, o de manera simplificada con $\la\cdot,\cdot\ra$ cuando el contexto no presente ambigüedades. Por ejemplo, cuando $X$ es un espacio de Hilbert su producto interno también se denotará con $\la\cdot,\cdot\ra$. Esto no presenta ninguna ambigüedad si tenemos en cuenta el \texttt{Teorema de representación de Riesz} que dice que para cada $f\in X^*,\,\exists!\, y\in X$ tal que $\|y\|_X=\|f\|_{X^*}$ y $f(x)=\la y,x\ra\,\forall x\in X$. Asumiremos en general, que $X$ es un espacio vectorial normado, y denotaremos por $\|\cdot\|$ a su respectiva norma (se utilizará un subíndice correspondiente cuando sea necesario aclarar la norma utilizada).


Comencemos con la topología. Será útil distinguir dos tipos de topologías, la fuerte, y la débil. Cuyas definiciones damos a continuación.

\begin{defi}[Topología fuerte]
Dados $\bx\in X$, $X$ espacio vectorial normado y $\varepsilon>0$ definimos $B[\bx,\varepsilon]=\{x\in X \mid \|x-\bx\|\leq\varepsilon\}$. Llamamos topología fuerte a la dada por estos $B$'s.

Diremos que $\{x_n\}\subset X$ converge a $\bx$ en la topología fuerte, o converge fuertemente (la denotaremos con $x_n\to \bx$), si dado $\varepsilon>0$ existe $N\in \N$ tal que $x_n\in B[\bx,\varepsilon]\,\forall n\geq N$.
\end{defi}
%  Diremos que si $\,x_n$ converge a $\bx$ en la topología fuerte, o converge fuertemente (lo denotaremos con $x_n\to \bx$), $\{x_n\}\subset X$ y $\bx\in X$, entonces, dado $c>0,\, \exists N_\epsilon \mbox{ tal que } \|x_n-\bx\|\leq\epsilon\,\,\forall n\geq N_\epsilon$, es decir, $x_n\in B(\bx,\epsilon)\,\,\forall n\geq N_\epsilon$.


\begin{defi}[Topología débil]
Dados $\bx\in X$, $\varepsilon>0$ y un conjunto finito $\{f_1,\ldots,f_k\}\subset X^*$ definimos $V(\bx;f_1,\ldots,f_k;\varepsilon)=\{x\in X \mid |\la f_i , x-\bx\ra|\leq\varepsilon\mbox{ para }i=1,\ldots,k\}$. Llamamos topología débil a la dada por estos $V$'s.

Diremos que $\{x_n\}\subset X$ converge a $\bx$ en la topología débil, o converge débilmente (la denotaremos con $x_n\deb \bx$), si dados $\varepsilon>0$ y $\{f_1,\ldots,f_l\}\subset X^*$ existe $N\in \N$ tal que $x_n\in V(\bx;f_1,\ldots,f_l;\varepsilon)\,\forall n\geq N$.
\end{defi}

La topología débil posee una propiedad fundamental, y es que cualquier sucesión convergente converge solamente a un punto. Esto se demuestra usando algunos resultados sobre convergencia en la topología débil. Podemos escribir esta propiedad formalmete de la siguiente manera. $$\mbox{Dada } \{x_n\}\subset X, \mbox{ si } x_n\deb x \mbox{ y }x_n\deb y, \mbox{ entonces } x=y.$$

En particular, cuando $X$ es de dimensión finita, la topología fuerte y la débil coinciden. Por lo tanto $\{x_n\}$ converge débilmente si y sólo si converge fuertemente.



%\textbf{Ejemplo}.
\begin{ejem}
Sea $S=\{x\in X\,|\, \|x\|=1\}$ y $B=\{x\in X\,|\,\|x\|\leq1\}$, entonces, la clausura de $S$ en la topología débil es $B$.
\end{ejem}

La siguiente proposición explica como funciona la convergencia en relación con estas dos topologías que usaremos (ver \cite[Proposition 3.5.]{brezis2010functional}).
\begin{prop}
 Sean $\{x_n\}\subset X$ y $\bx\in X$, entonces
 \begin{enumerate}
 \item $x_n\deb \bx$ si y sólo si $\la f,x_n\ra\to\la f,\bx\ra\,\forall f\in X^*$.
 \item $x_n\to \bx$ implica $x_n\deb \bx$.
 \item $x_n\deb\bx$ implica que $\{\|x_n\|\}$ es una sucesión acotada y $\|\bx\|\leq\varliminf_{n\to\infty}\|x_n\|$.
 \item Si $x_n\deb\bx$ y $f_n\to f$ entonces $\la f_n,x_n\ra\to\la f,\bx\ra$. 
 \end{enumerate}
\end{prop}

Introduciremos ahora, algunos resultados que requieren la definición de conjunto convexo.
Los conjuntos y funciones convexas, poseen propiedades muy importantes y de utilidad en la optimización.

Diremos que un conjunto $D\subset X$ es convexo si para cualquier par de puntos $x,y\in D$ y $\alpha\in[0,1]$ se tiene que $$\alpha x + (1-\alpha)y\in D.$$
Y una función $f:D\arr\R$ se dice convexa si para cualquier par de puntos $x\, ,\,y\in D$ con $y\neq x$ y para cualquier $\alpha\in (0,1)$, se tiene que $$f(\alpha x+(1-\alpha)y)\leq\alpha f(x)+(1-\alpha)f(y).$$

Ahora podemos enunciar el siguiente teorema, que relaciona la convexidad, con las topologías dadas anteriormente (ver \cite[Theorem 3.7.]{brezis2010functional}).

\begin{teo}
Si $C\subset X$ es convexo, entonces, $C$ es cerrado en la topología débil si, y sólo si, $C$ es cerrado en la topología fuerte.
\end{teo}

Para trabajar con la topología débil, necesitamos definir ciertos conceptos topológicos, como compacidad, clausura, y el significado de ser un conjunto cerrado, cuando la topología que estamos usando es la débil.

\begin{defi}
 $C$ es débilmente secuencialmente cerrado si para una sucesión $\{x_n\}\subset C$ con $x_n\deb~\bx\in X$ vale que $\bx\in C$.
\end{defi}
\begin{defi}
 La clausura débil secuencial se define como
 \[
\mathrm{cl}_W(C)=\{x\in X\mid \exists\{x_n\}\subset C\mbox{ tal que }x_n\deb x\}.  
 \]

\end{defi}
\begin{defi}
 $C$ es débilmente secuencialmente compacto si toda sucesión $\{x_n\}\subset C$ tiene una subsucesión débilmente convergente a algún $\bx\in C$.
\end{defi}

Note que si C es cerrado en la topología débil, entonces es débilmente secuencialmente cerrado.

Para optimizar en dimensión infinita, le pediremos a nuestros espacios de salida que sean \emph{reflexivos}, cuya definición damos a continuación.

\begin{defi}
 Sea $X$ un espacio de Banach, y $J:X\arr X^{**}$ la inyección canónica, i.e., $\la J(x),f\ra_{X^{**},X^*}=\la f,x\ra_{X^*,X}$ $\forall x\in X$, $\forall f\in X^*$. Diremos que $X$ es reflexivo si $J(X)=X^{**}$.
\end{defi}

En el caso de $X$ reflexivo, usaremos $\la x,f\ra$ $\forall x\in X$ $\forall f\in X^*$ en vez de $\la J(x),f\ra_{X^{**},X^*}$.

Lo que queremos ver ahora, es una manera de corroborar que un subespacio $D\subset X$ es débilmente secuencialmente compacto, sin usar la definición vista anteriormente, pues no es fácil de verificar. Esto se puede hacer en espacios de Banach reflexivos usando el siguiente resultado (ver \cite[Theorem 3.17.]{brezis2010functional}.
\begin{teo}[Kakutani]\label{kakutani}
 Sea $X$ un espacio de Banach. Entonces $X$ es reflexivo si y sólo si
 \[
  B=\{x\in X\,|\,\|x\|\leq1\}
 \]
es compacto en la topología débil.
\end{teo}

Como consecuencia de este postulado podemos formular las siguientes proposiciones (la primera se puede ver en \cite[Theorem 3.18.]{brezis2010functional}).

\begin{prop}
Sea $X$ un espacio de Banach reflexivo y $\{x_n\}\subset X$ una sucesión acotada. Entonces $\exists$ una subsucesión $\{x_{n_k}\}$ débilmente convergente.
\end{prop}


\begin{cor}
Sea $X$ un espacio de Banach reflexivo. Si $D\subset X$ es acotado, convexo, cerrado (da lo mismo la topología fuerte o débil por teorema anterior), y no vacío, entonces, $D$ es débilmente secuencialmente compacto.
\end{cor}

Por lo tanto, para obtener la compacidad débil, sólo hace falta tener un conjunto acotado, convexo, cerrado y no vacío en un espacio de Banach reflexivo.

Ya hemos caracterizado el conjunto donde queremos que esté definida nuestra función. Con el objetivo de probar la existencia de un minimizador global, necesitamos también caracterizar la función a minimizar. Para lo cual damos estas definiciones.

\begin{defi}
 Diremos que $F:X\arr Y$ es débilmente secuencialmente continua si $\forall\, \{x_n\}\subset X$ sucesión, con $x_n\deb x\in X$, vale que $F(x_n)\deb F(x)$.
\end{defi}

\begin{defi}
 Sea $D\subset X$, decimos que $f:D\arr\R$ es débilmente semicontinua inferiormente si $\forall\, \{x_n\}\subset X$ sucesión, tal que $x_n\deb x\in D$,  vale que $\varliminf_{n \to \infty} f(x_n)\geq f(x)$.
\end{defi}

Ahora, de la misma manera que encontramos una forma de saber si un conjunto es débilmente secuencialmente compacto, con el siguiente resultado, podemos saber si una función es débilmente semicontinua inferiormente, ya que también es una tarea compleja hacerlo desde la definición (ver \cite[Corollary 3.9.]{brezis2010functional}.

\begin{teo}\label{convycont}
Si $D\subset X$ es un conjunto cerrado y convexo, y $f:D\arr\R$ es una función convexa y continua, entonces es débilmente semicontinua inferiormente.
\end{teo}

Finalmente el teorema de existencia es
\begin{teo}\label{globalsol}
Sea $D\subset X$ y $f:D\arr\R$. Si $D$ es débilmente secuencialmente compacto y $f$ es débilmente semicontinua inferiormente, entonces $\exists\, \bx\in D$ tal que $f(\bx)\leq f (x)\,\forall x\in D$.
\end{teo}
\begin{dem}
Sea $\{x_n\}\subset D$ tal que $f(x_n)\leq \mbox{inf}_{x\in D}f(x)+\frac{1}{n}$. Como $D$ es débilmente secuencialmente compacto, existe una subsucesión $\{x_{n_k}\}\subset D$ y $\bx\in D$ tal que $x_{n_k}\deb\bx$. Luego, por ser $f$ débilmente semicontinua inferiormente, se obtiene que $\mbox{inf}_{x\in D}f(x)\geq\varliminf_{k\to\infty}f(x_{n_k})\geq f(\bx)$.
\end{dem}

Ahora, volviendo a nuestro problema \eqref{problema}, nos gustaría ver que nustra función objetivo $f(c,a,p,v_0)=\int_{0}^{1} \frac{p}{\sqrt{2p \int_0^s a(r)dr + v_0^2}}ds$ es débilmente semicontinua inferiormente. En principio, podemos afirmar que es continua, ya que las funciones de movimiento son continuas, pero no necesariamente es convexa, lo que imposibilita utilizar el Teorema \eqref{convycont} para demostrar que es débilmente semicontinua inferiormente. De la misma manera el conjunto donde minimizamos la función (teniendo en cuenta las restricciones) no es convexo.

Ya que la función es continua, puede parecer sensato ver que este conjunto donde minimizamos es compacto, ya que si es así, nos aseguramos la existencia de un mínimo. El problema es que verificar compacidad en espacios de funciones resulta demasiado complicado. El enfoque podría basarse en que si las funciones no están acotadas, entonces esto generaría un absurdo, por lo que debemos utilizar la topología débil (por \eqref{kakutani}), pero la convergencia en topología débil no tiene las buenas propiedades de la fuerte, con lo que no logramos llegar al absurdo.

Sin embargo, si demostraremos la existencia del problema discretizado, utilizando la convergencia en la topología fuerte, esto se verá más adelante en la Sección \eqref{existsoldiscretizado}. 

%, y esto es porque, dada una función posición (o trayectoria), si en un instante de tiempo se encuentra en una posición $x_0$, en el siguiente instante no puede encontrarse en otra posición dando un salto (es decir que sería un punto de discontinuidad) pues el movimiento descripto por la física del tamaño de partículas que estamos manejando, no puede tener saltos de discontinuidad. El mismo argumento podemos utilizar para las demás funciones de movimiento.

%Por lo tanto $f$ es continua.


{\bf FALTA VINCULAR TODA ESTA SECCION AL RESTO DEL TRABAJO}
% Volviendo entonces al planteo del problema \eqref{problema} podemos asegurar que nuestra función objetivo $f:H^2([0,1])^2 \times L^2([0,1])\times \R_+\times \R_+\arr\R$ con $f(c,a,p,v_0)=\displaystyle\int_{0}^{1} \frac{p}{\sqrt{2p \int_0^s a(r)dr + v_0^2}}ds$ es continua, pero no es necesariamente convexa, por lo que no podemos usar el teorema que antes vimos para decir que nuestra función es débilmente semicontinua inferiormente. Ya que verificar esta propiedad, al igual que verificar que el conjunto donde estamos minimizando ($H^2([0,1])^2 \times L^2([0,1])\times \R_+\times \R_+$ en este caso) es débilmente secuencialmente compacto es muy complejo, lo dejaremos de lado y asumiremos que el problema tiene solución, lo que resulta (en términos de la física del problema) convincente, ya que podemos esperar que exista una trayectoria óptima.


\subsection{Existencia de multiplicadores de Lagrange}
Queremos resolver el problema \eqref{problema}, una forma de hacerlo, es ver que existen los multiplicadores de Lagrange asociados al problema, y así, con métodos computacionales, podremos encontrar una solución.


Para comenzar, al igual que para los teoremas de optimización en dimensión finita debemos ser capaces de derivar, necesitamos definir diferenciabilidad en dimensión infinita.
% Poner definicion de ``o'' (o pequeño) en Notación
% poner \mathcal{L}(X,Y) en Notación
% Poner la Suma de Minkowski en Notación
\begin{defi} Sean $X$ e $Y$ dos espacios vectoriales normados (EVN's), entonces diremos que la función $F:X\arr Y$ es \textbf{Fréchet diferenciable} (o simplemente, diferenciable) en $\bx$ si existe $L\in\mathcal{L}(X,Y)$ tal que
$$F(x)=F(\bx)+L(x-\bx)+o(\| x-\bx\|).$$
En este caso $L$ es llamada la derivada de Fréchet de $F$ en $\bx$ y se la denota por $F'(\bx)$.
\end{defi} 
En el caso de que $Y=\R$ entonces $F'(\bx)\in X^*$ y por lo tanto la diferenciabilidad queda caracterizada por
$$F(x)=F(\bx)+\la F'(\bx),x-\bx\ra+o(\|x-\bx\|).$$


 Necesitaremos también unos conjuntos particulares,  el cono normal y el cono tangencial que serán usados en varios teoremas Pero antes, definamos \texttt{cono}.

\begin{defi}
Sea $X$ un espacio vectorial, y $K\subset X$ no vacío. Decimos que $K\subset X$ es un \textbf{cono} si se cumple que $$\alpha K=\{\alpha x\mid x\in K\}\subset K\quad\forall\,\alpha\geq 0.$$ 
\end{defi} 

Como veremos más adelante en el teorema de los multiplicadores de Lagrange, pediremos que nuestro cono sea convexo, y una manera fácil de verificar esta propiedad es con la siguiente proposición.

\begin{prop}
 Un cono $K$ es convexo si, y sólo si, $K+K\subset K$.
\end{prop}

\begin{ejem}
El conjunto $K=\{v\in\Rn \mid v_i=0,\,i=1,\ldots,l;\,v_j\leq0,\, j=l+1,\ldots,n\}$ es un cono convexo, pues, es un cono ya que, si $\alpha\geq0,\,\alpha v_i=0$ y $\alpha v_j\leq0$ y por lo tanto $\alpha K\subset K$; y es convexo porque, si $v,w\in K$, entonces $v+w\in K$ pues si $v_i=0$ y $w_i=0,\,v_i+w_i=0$ y si $v_j\leq0$ y $w_j\leq0,\, v_j+w_j\leq0$.
 
\end{ejem}


Los conos particulares que nombramos antes se definen así (note que, efectivamente son conos en el sentido de la definición dada anteriormente).

\begin{defi}
Sea $X$ un espacio de Banach, y $S\subset X$. Definimos el \textbf{cono tangente} a $S$ en $\bx$ como
$$\cT_{S} (\bx)=\{d\in X\mid\exists\,\{(t_{k},d^k)\}\subset\R_+\times X\;\mathrm{tal\; que}\;\bx+t_{k}d^k\in S\;\mathrm{y}\;t_{k}\arr 0^+,\;d^k\arr d\},$$ si $\bx\in S$ y $\cT_{S}(\bx)=\emptyset$ si $\bx\notin S$.
\end{defi}

\begin{defi}
Sea $X$ un espacio de Banach, y $S\subset X$. Definimos el \textbf{cono normal} como 
$$\cN_{S}(\bx)=\left\{v\in X^*\;\left|\begin{array}{c}
\forall\left\{x^k\right\}\subset S\backslash\{\bx\}\mbox{ con } x^k\arr\bx,\\
\mbox{vale que }\displaystyle\limsup_{k\arr\infty}\frac{\langle v,x^k-\bx\rangle}{\|x^k-\bx\|}\leq 0\end{array}\right.\right\},$$
si $\bx\in S$ y $\mathcal{N}_{S}(\bx)=\emptyset$ si $\bx\notin S$.
\end{defi}

Avanzando hacia el teorema de los multiplicadores de Langrange, necesitamos la noción del cono dual de un cono y una útil propiedad de los conos definidos anteriormente.

\begin{defi}
 Definimos el cono dual de un cono  $K\subset X$ como
$$K^*=\{f\in X^*\mid\langle f,x\rangle\leq 0\;\forall\,x\in K\}.$$
\end{defi}

Cuando $X$ es reflexivo, el bidual de un cono $K$ se define como $K^{**}=(K^*)^*$. Una propiedad del bidual que usaremos luego es la siguiente.
\begin{prop}\label{conobidual}
 Sea $X$ reflexivo y $K\subset X$ un cono convexo. Entonces $K^{**}=\mathrm{cl}(K)$. En particular si $K$ es cerrado, entonces $K^{**}=K$.
\end{prop}
La demostración de esta Proposición puede verse en \cite[Proposition 2.40]{bonnans2013perturbation}.

\begin{teo}
Para todo $\bx\in S$ se cumple que $(\cT_{S}(\bx))^*=\cN_{S}(\bx)$.
\end{teo} 
\begin{dem}
Veamos primero la inclusión $\cT_S(\bx)^*\subset \cN_S(\bx)$. Sea $v\in X^*$, si $v\notin\cN_S(\bx)$ entonces (sin pérdida de generalidad) existe una sucesión $\{x_n\}\subset S\setminus\{\bx\}$ con $x_n\to \bx$ y $\lim_{n\to\infty}\frac{\la v, x_n-\bx\ra}{\|x_n-\bx\|}>0$. Definimos $$d_n=\frac{x_n-\bx}{\|x_n-\bx\|},$$ como $\{d_n\}$ es acotada (pues $\|d_n\|=1\,\forall n$) y $X$ es reflexivo, existe una subsucesión $d_{n_j}\deb d$, entonces $\la v , d_{n_j}\ra\to\la v , d\ra>0$. Por otro lado, $\bx+\|x_n-\bx\|d_n=x_n\in S$ por lo que $d\in \cT_S(\bx)$. Ahora $\la v , d\ra>0$ con $d\in\cT_S(\bx)$ y entonces $v\notin\cT_S(\bx)^*$.

Veamos ahora la inlusión $\cN_S(\bx)\subset\cT_S(\bx)^*$. Tomemos $v\in\cN_S(\bx)$. Si $d=0$ entonces $\la v,d\ra\leq0$. Si $d\in\cT_S(\bx)\setminus\{0\}$, existe $t_n\to 0^+$ y $d_n\to d$ tal que $x_n:=\bx+t_nd_n\in S$. Como $t_nd_n\neq0$ y $t_nd_n\to 0$ para $n$ suficientemente grande, tenemos que $\{x_n\}\subset S\setminus\{x\}$ y $x_n\to\bx$. Por lo tanto $$0\geq\varlimsup_{n\to\infty}\frac{\la v,x_n-\bx\ra}{\|x_n-\bx\|}=\frac{\la v,d_n\ra}{\|d_n\|}=\frac{\la v,d\ra}{\|d\|}$$ entonces $\la v,d\ra\leq0$ lo que implica que $v\in \cT_S(\bx)$.
\end{dem}

Con esto, ya podemos enunciar el teorema de optimalidad de primer orden, en versión infinita.

\textbf{Nota}. A partir de este momento utilizaremos la terminología local y global, la cual hace referencia al cumplimiento de cierta propiedad en un entorno dado por la topología correspondiente, o en todo el espacio con el que se trabaja; respectivamente. 

\begin{teo}[Condición necesaria de optimalidad en forma primal]
Sea $S\subset X$ y $f:S\arr\R$, Fréchet diferenciable en $\bx\in S$. Si $\bx$ es una solución local de \[
\begin{array}{rl}
\displaystyle \mathop{\mathrm{minimizar}}_{x\in S} & \displaystyle f(x)
\end{array}
\]   
Entonces $$\langle f'(\bx) , d\rangle_{X^*,X}\geq 0 \quad\forall d\in \cT_S(\bx).$$
Equivalentemente,
$$ 0\in f'(\bx)+ \cN_{S}(\bx).$$
\end{teo}
\begin{dem}
Tomemos $d\in\cT_S(\bx)$, entonces existe $t_n\to 0^+$ y $d_n\to d$ tal que $\bx+t_nd_n\in S$. Como $\bx$ es mínimo local y $\bx+t_nd_n\to\bx$ para $n$ suficientemente grande, tenemos que $$f(\bx)\leq f(\bx+t_nd_n)=f(\bx)+t_n\la f'(\bx),d_n\ra+o(\|t_nd_n\|).$$ Ahora por $t_n>0$ obtenemos $$0\leq\la f'(\bx),d_n\ra+\frac{o(\|t_nd_n\|)}{\|t_n d_n\|}\|d_n\|\to\la f'(\bx),d\ra.$$
\end{dem}

Se puede observar que la definición de los conos normal y tangencial son bastante complejas, por lo que resulta difícil trabajar con ellas y verificar que algún elemento pertenece o no a determinado cono. La siguiente proposición simplifica la definición simpre y cuando el conjunto sea convexo.



\begin{prop}\label{conoconv}
Sea $S\subset X$ un conjunto convexo, $X$ un espacio de Banach reflexivo, $\bx\in S$. Entonces vale que:
\begin{enumerate}
\item $\cT_{S}(\bx)=\mathrm{cl}(\mathrm{cone}(S-\{\bx\})$.
\item $\cN_{S}(\bx)=\{v\in X^*\mid\langle v,x-\bx\rangle\leq 0\; \forall\,x\in S\}$.
\end{enumerate}
Siendo $\mathrm{cone}(S)=\{\alpha x\mid \alpha\geq0,\,x\in S\}$.
\end{prop} 
\begin{dem}
\begin{enumerate}
\item Veamos la inclusión $\mathrm{cl}(\mathrm{cone}(S-\{\bx\})\subset\cT_{S}(\bx)$. Sea $d\in K:=\mathrm{cone}(S-\{\bx\})$, entonces existe $\alpha\geq0$ y $x\in S$ tal que $d=\alpha(x-\bx)$. Si $d=0,\, 0\in \cT_S(\bx)$ y listo.
Si $d\neq0$ entonces $\alpha>0$. Tomo $t_n\to0^+$ y $d_n=d$. Para $n$ suficientemente grande $$\alpha t_n\in (0,1)\mbox{ y }\bx+t_nd=(1-\alpha t_n)\bx+\alpha t_n(\bx+\frac{1}{\alpha}d)\in S,$$ por lo que $d\in S$. Como $\cT_S(\bx)$ es cerrado (por ser un conjunto de puntos de acumulación) y $K\subset\cT_S(\bx)$ obtenemos $$\mathrm{cl}(K)\in\mathrm{cl}(\cT_S(\bx))=\cT_S(\bx).$$

Veamos ahora la inclusión $\cT_{S}(\bx)\subset\mathrm{cl}(\mathrm{cone}(S-\{\bx\})$. Tomo $d\in\cT_(\bx)$ entonces existe $t_n\to0^+$ y $d_n\to d$ tal que $\bx+t_nd_n\in S$, luego $$d_n\in\frac{1}{t_n}(S-\{\bx\})\subset K\subset\mathrm{cl}(K),$$ y por lo tanto $d\in \mathrm{cl}(K)$.
\item Veamos la inclusión $\{v\in X^*\mid\langle v,x-\bx\rangle\leq 0\; \forall\,x\in S\}\subset\cN_{S}(\bx)$. Sea $v\in Q:=\{v\in X^*\mid \la v,x-\bx\ra\leq0\,\forall x\in S\}$ y tomo $d\in\cT_(\bx)$, entonces existe $t_n\to 0^+$ y $d_n\to d$ tal que $x_n:=\bx+t_nd_n\in S$. Ahora $$0\geq \la v , x_n-\bx\ra=t_n\la v,d_n\ra$$ y entonces $\la v ,d_n\ra\leq0$ (pues $t_n>0)$. Luego $\la v,d\ra\leq0$ entonces $v\in\cT_(\bx)^*=\cN_S(\bx)$.

Ahora veamos $\cN_{S}(\bx)\subset\{v\in X^*\mid\langle v,x-\bx\rangle\leq 0\;\forall\,x\in S\}$. Sea $v\in\cN_S(\bx)=\cT_S(\bx)^*$. Tomo $x\in S$, entonces $x-\bx\in K\subset\cT_S(\bx)$. Luego $$\la v ,x-\bx\ra\leq0$$ por lo que $v\in Q$.
\end{enumerate}
\end{dem}

\textbf{Importante}.
Consideremos el problema
\[\tag{P}
\begin{array}{rl}
\displaystyle \mathop{\mathrm{minimizar}}_{x} & \displaystyle f(x) \\
\mathrm{sujeto \ a} & G(x)\in K,\, x\in C.
\end{array}
\]
donde
\begin{itemize}
\item$f:X\arr\R$ y $G:X\arr Z$ continuamente Fréchet diferenciables.

\item $X$ Banach reflexivo y $Z$ Banach,

\item $C\subset X$ convexo cerrado y $K\subset Z$ cono convexo cerrado.
\end{itemize}


El conjunto factible de (P) es $$S=\{x\in C | G(x)\in K\}$$

Note que el problema (P) tiene la misma estructura que el problema que planteamos en la sección anterior, y queremos resolver. Con un poco más de teoría ya seremos capaces de enunciar el teorema de los multiplicadores de Lagrange en dimensión infinita.

\begin{defi}
 Llamaremos cono de linealizaciones asociado al problema (P) en el punto $\bx \in S$ al conjunto $$L_S(\bx)=\{d\in\cT_C(\bx) | G'(\bx)d\in \cT_K(G(\bx))\}.$$
 (Note que este conjunto también es efectivamente un cono).
\end{defi} 

\begin{prop}
Se $X$ Banach reflexivo, entonces $L_S(\bx)^* = \{G'(\bx)^*\mu + \omega\mid\mu\in\cN_K(G(\bx)),\, \omega\in\cN_C(\bx)\}$
\end{prop} 
\begin{dem}
Sea $B=\{G'(\bx)^*\mu+\omega\mid\mu\in\cN_K(G(\bx)),\,\omega\in \cN_C(\bx)\}$. Comencemos viendo $$\{G'(\bx)^*\mu + \omega\mid\mu\in\cN_K(G(\bx)),\, \omega\in\cN_C(\bx)\}\subset L_S(\bx)^*.$$ Sea $v\in B$, entonces existe $\mu\in\cN_K(G(\bx)),\,\omega\in\cN_C(\bx)$ tal que $v=G'(\bx)\mu+\omega.$ Sea $d\in L_S(\bx)$, entonces $$\la v,d\ra=\la G'(\bx)\mu+\omega,d\ra=\la \mu,G'(\bx)d\ra+\la\omega,d\ra\leq0$$ pues los dos miembros son menores o iguales que cero. Por lo tanto $v\in L_S(\bx)^*$.

Veamos ahora $$L_S(\bx)^* \subset \{G'(\bx)^*\mu + \omega\mid\mu\in\cN_K(G(\bx)),\, \omega\in\cN_C(\bx)\}.$$ Sea $d\in B^*$, entonces $$0\geq\la G'(\bx)^*\mu+\omega,d\ra=\la\mu,G'(\bx)d\ra+\la\omega,d\ra\,\forall\mu\in\cN_K(G(\bx)),\,\forall \omega\in\cN_C(\bx).$$ Haciendo $\omega=0$ y luego $\mu=0$, obtengo $0\geq\la\mu,G'(\bx)d\ra\,\forall\mu\in\cN_K(G(\bx))$, entonces $$G'(\bx)d\in\cN_K(G(\bx))^*=\cT_K(G(\bx))\mbox{ y }0\geq\la\omega,d\ra\,\forall\omega\in\cN_C(\bx).$$ Luego $d\in\cN_C(\bx)^*=\cT_C(\bx)$ y así $d\in L_S(\bx)$. Por lo tanto $B^*\subset L_S(\bx)$. Como $B$ es un cono convexo cerrado, usando la Proposición \ref{conobidual} se obtiene que $L_S(\bx)^*\subset B^{**}=B$.
\end{dem}

La siguiente definición es fundamental para nuestro teorema de multiplicadores de Lagrange, pues explica una condición suficiente para garantizar la tesis del teorema.
En la literatura existen muchas condiciones suficientes como ésta, y se investiga la existencia de alguna necesaria y suficiente, pues hay casos en los que existen multiplicadores de Lagrange a pesar de que no se cumplan ninguna de estas condiciones suficientes.
\begin{defi}
 Diremos que la condición de regularidad de Robinson vale en $\bx\in S$ si $G$ es Fréchet diferenciable en $\bx$ y 
 \begin{equation}\label{robinson}
0\in \mathrm{int}\left(G(\bx) + G'(\bx)(C-\bx)-K\right)=\mathrm{int}(\{G(\bx) + G'(\bx)(x-\bx)-v \mid x\in C , v\in K\}).
 \end{equation}
\end{defi} 

El siguiente resultado relaciona el cono de linealizaciones con el cono tangente y su demostración puede encontrarse en \cite[Corollary 2.91]{bonnans2013perturbation}.
\begin{prop}
 Si la condición de regularidad de Robinson (\ref{robinson}) vale en $\bx \in S$, entonces $\cT_S(\bx) = L_S(\bx)$.
\end{prop}

Finalmente, el teorema de los multiplicadores de Lagrange es el siguiente.

\begin{teo}\label{lagrinf}
Sea $\bx$ solución local de (P), donde $f:X\arr\R$ y $G:X\arr Z$ son continuamente Fréchet diferenciables en $\bx$, con $X$ Banach reflexivo, $C\subset X$ convexo cerrado, $K\subset Z$ cono convexo cerrado.

Si vale la condición de Robinson (\ref{robinson}) en $\bx$, entonces $\exists \mu\in Z^*$ tal que 
\begin{equation}\label{condOptim}
\begin{array}{c}
         -(f'(\bx) + G'(\bx)^*\mu)\in \cN_C(\bx)\\
          G(\bx)\in K,\quad \mu\in K^* ,\quad \langle\mu , G(\bx)\rangle = 0
         \end{array}
\end{equation}
\end{teo} 
\begin{dem}
Como $\bx$ es solución local de $\displaystyle \mathop{\mathrm{minimizar}}_{x\in S}\displaystyle f(x)$, vale que $$\la f'(\bx),d\ra\geq0\,\forall d\in\cT_S(\bx).$$ Como también vale la condición de Robinson, tenemos que $$\la f'(\bx),d\ra\geq0\,\forall d\in L_S(\bx),$$ y entonces $-f'(\bx)=G'(\bx)^*\mu+\omega$, lo que implica que $$-(f'(\bx)+G'(\bx)^*\mu)=\omega\in\cN_C(\bx)\mbox{ y }\mu\in\cN_K(G(\bx)),$$ por lo que $\mu\in K^*$ y $\la \mu,G(\bx)\ra=0$.
\end{dem}

\begin{defi}
 Diremos que $L:X\times Z^*\arr \R$ tal que $L(x,\mu)= f(x)+\langle\mu , G(x)\rangle$ es la función lagrangiana asociada al problema (P), y diremos que $\mu\in Z^*$ es un multiplicador de Lagrange asociado a (P) en $\bx$ si 
 $$-\frac{\partial L}{\partial x}(\bx,\mu)\in\cN_C(\bx),\ G(\bx)\in K,\ \mu\in K^* ,\  \langle\mu , G(\bx)\rangle=0.$$ 
\end{defi}

Recordando nuestro problema \eqref{problema} nos gustaría comprobar que vale la condición de Robinson (\ref{robinson}), lo cual nos aseguraría la existencia de multiplicadores de Lagrange, pero como es común en los problemas de matemática aplicada, no es posible verificar esta condición debido a la complejidad de los espacios que se manejan (note que para nuestro problema $S$ es el conjunto de elementos $x\in H^2([0,1])^2 \times L^2([0,1])\times \R_+\times \R_+$ tal que nuestras restricciones evaluadas en ese $x$ se cumplan). Por lo cual, asumiremos que existen multiplicadores de Lagrange, con el objetivo de resolver el problema con métodos numéricos que utilizan el hecho de la existencia de estos multiplicadores.

\section{Dimensión finita}

Comenzamos esta sección que pretende dar los teoremas de existencia de multiplicadores de Lagrange, pero ahora en dimensión finita, pues como se ha mencionado anteriormente, su existencia es necesaria para aplicar muchos de los métodos computacionales que tienen como objetivo resolver un problema de optimización. El siguiente resultado puede obtenerse de lo visto para el caso en dimensión infinita.

\begin{cor}\label{solglobalfinita}
 Sea $f:\R^n\to\R$ continua, $D\subset\R^n$ cerrado y $x_0\in D$. Si el conjunto $\{x\in D \mid f(x)\leq f(x_0)\}$ es compacto y no vacío, entonces existe $\bar{x}\in D$ tal que $f(\bar{x})\leq f(x)$ para todo $x\in D$.
\end{cor}
\begin{dem}
 Sea $C=\{x\in D \mid f(x)\leq f(x_0)\}$. Como estamos en un espacio de dimensión finita, la topología débil coincide con la topología fuerte. Por lo tanto, como $C$ es compacto entonces es débilmente secuencialmente compacto y como $f$ es continua entonces $f$ es débilmente semicontinua inferiormente. Luego, el Teorema \ref{globalsol} nos garantiza la existencia de $\bx\in C$ tal que $f(\bx)\leq f(x)$ para todo $x\in C$.
 
 Claramente, si $x\in D\setminus C$ entonces $f(x)>f(x_0)\geq f(\bx)$ pues $x_0\in C$. Por lo tanto $f(\bx)\leq f(x)$ para todo $x\in D$.
\end{dem}


El caso de la existencia de soluciones en dimensión finita, es de igual manera compleja de demostrar matemáticamente, ya que ni siquiera discretizada, nuestra función objetivo es convexa, por lo que no podemos asegurar que sea débilmente semicontinua inferiormente (como vimos el la sección \ref{existenciadesoluciones}). Por lo que el siguiente paso es continuar con la teoría finita de los multiplicadores de Lagrange.

Consideremos el problema (P) discretizado, es decir,
\[\tag{$\tilde{P}$}
\left\{\begin{array}{rl}
\displaystyle \mathop{\mathrm{minimizar}}_{x} & \displaystyle f(x) \\
\mathrm{sujeto \ a} & G_i(x)=0, \, i\in \cE,\\
                    & G_i(x)\leq 0, \, i\in \cI,
\end{array}\right.
\]
donde $f:\R^n\arr\R$ y $G:\R^n\arr\R^m$ son diferenciables y $$\cE\cup \cI=\{1,\ldots,m\}\mbox{ con } \cE\cap \cI=\emptyset.$$
Denotaremos por $m_\cE$ y $m_\cI$ la cardinalidad de los conjuntos $\cE$ y $\cI$, respectivamente. Además, para $v\in\R^m$, denotamos por $v_\cE$ a la $m_E$-upla conformada por las componentes $v_i$ de $v$, con índices $i\in \cE$; análogamente para $v_\cI$. Para $A\in\R^{m\times r}$ denotamos por $A_\cE$ a la submatriz de $m_E$ filas y $r$ columnas conformada por las filas con índices $i\in \cE$; análogamente para $A_\cI$.


Note que este problema es un caso particular del problema (P) considerando
\begin{itemize}
    \item $X=\Rn,$
    \item $C=\Rn,$
    \item $Z=\R^m,$
    \item $K=\{v\in\R^m\mid v_i=0,\,i\in\cE;\ v_i\leq 0,\,i\in \cI\}$.
\end{itemize}
Tomaremos en cuenta a $\Rn$ dotado del producto interno euclídeo, i.e., $\la x,y\ra=x^ty$ $\forall\, x,y\in\Rn$. 
Por lo tanto el conjunto factible de este problema puede escribise como
\[
 S=\{x\in\Rn \mid G_i(x)=0,\; i\in\cE;\, G_i(x)\leq0,\; i\in\cI\}.
\]


\begin{defi}
Para $x\in S$, el conjunto de restricciones de desigualdad activas en $x$ es $\cA(x)=\{j\in\cI \mid G_j(x)=0\}$.
\end{defi}

En este contexto, la condición de Robinson puede entenderse de la siguiente manera.

\begin{prop}
La condición de Robinson (\ref{robinson}) vale en $\bx\in S$ si y solo si se cumplen
\begin{enumerate}[a)]
\item $\{G_i'(\bx)\}_{i\in\cE}$ son linealmente independientes,
\item existe $d\in\Rn$ tal que $\la G_i'(\bx),d\ra=0$ para $i\in\cE$ y $\la G_i'(\bx),d\ra<0$ para $i\in\cA(\bx)$.
\end{enumerate}
\end{prop}
\begin{dem}
Veamos primero que si vale la condición de Robinson en $\bx$ se cumplen (a) y (b).
\begin{enumerate}[a)]
\item Porque vale la condición de Robinson, tenemos que $\forall q\in B[0,\varepsilon],\,\exists d \in\Rn,\,v\in K$ tal que $q=G(\bx)+ G'(\bx)d-v$. Entonces como $G(\bx),v\in K$ si $i\in\cE$ vale $G_i(\bx)=0=v_i$ por lo que $q_\cE=G_\cE'(\bx)d$ y esto ocurre si, y solo si $G_\cE'(\bx)$ es de rango completo. Entonces $G_i'(\bx)$ son linealmente independientes para $i\in\cE$.
\item Consideremos $q=(q_\cE,q_\cI)\in B[0,\varepsilon]$ con $q_\cE=0,\,q_\cI\leq0$. Tomamos como antes $q=G(\bx)+\la G'(\bx),d\ra-v$ entonces $0=\la G_i'(\bx),d\ra$ si $i\in \cE$.
Si $i\in \cA(\bx)\subset \cI$, $\la G_i'(\bx),d\ra=q_i-G_i(\bx)+v_i\leq q_i<0$ pues $i\in \cI$.
\end{enumerate}

Veamos ahora que si valen (a) y (b), se cumple la condición de Robinson.

Por (a) sabemos que
$$C_G:=\|G_\cE'(\bx)^t(G_\cE'(\bx)G_\cE'(\bx)^t)^{-1}\|<+\infty.$$
Sea $d\in\Rn$ dado por la condición (b) y definimos
$$\beta=\mathrm{max}\left\{\mathrm{max}\{\la G_j'(\bx),d\ra\mid j\in \cI\setminus\cA(\bx)\},1\right\},$$
$$\gamma=\mathrm{max}\{\la G_j'(\bx),d\ra\mid j\in \cA(\bx)\},$$
$$\delta=\mathrm{max}\{G_j(\bx)\mid j\in \cI\setminus\cA(\bx)\},$$
$$\eta=\mathrm{max}\{\|G_j'(\bx)\|\mid j\in \cI\}.$$
Ahora, tomo
$$\alpha\in \left(0,\frac{-\delta}{2(\beta-\frac{\gamma}{2})}\right], \qquad \varepsilon\in\left(0,\frac{-\alpha\gamma}{2(\eta C_G+1)}\right].$$

Para $q\in B[0,\varepsilon]$, definimos $\tilde{d}=G_\cE'(\bx)^t(G_\cE'(\bx)G_\cE'(\bx)^t)^{-1}q_\cE$ y $v=G(\bx)+ G'(\bx)(\tilde{d}+\alpha d)-q$. Veamos que $v\in K$.

Como $q_\cE=G_\cE'(\bx)\tilde{d}$, entonces para $i\in \cE$,
\begin{eqnarray*}
v_i&=&G_i(\bx)+ \la G_i'(\bx),\tilde{d}\ra +\alpha \la G_i'(\bx),d\ra -q_i\\
&=& \la G_i'(\bx),\tilde{d}\ra -q_i = 0,
\end{eqnarray*}
donde usamos que $d$ cumple (b). Note que
$$\|\tilde{d}\|\leq C_G \|q_\cE\|\leq C_G\varepsilon.$$
Si $j\in \cI\setminus\cA(\bx)$,
\begin{eqnarray*}
v_j&=&G_j(\bx)+\la G_j'(\bx),\tilde{d}\ra + \alpha\la G_j'(\bx),d\ra-q_j\\
&\leq& \delta+\eta C_G\varepsilon+\alpha\beta+\varepsilon\\
&=&\delta + \varepsilon(\eta C_G+1)+\alpha\beta\\
&\leq& \delta + \alpha(\beta-\frac{\gamma}{2})\\
&\leq& \delta -\frac{\delta}{2} =\frac{\delta}{2}<0.
\end{eqnarray*}
Si $j\in \cA(\bx)$,
\begin{eqnarray*}
v_j&=&\la G_j'(\bx),\tilde{d}\ra + \alpha\la G_j'(\bx),d\ra-q_j\\
&\leq& \eta C_G\varepsilon+\alpha\gamma+\varepsilon\\
&=& \varepsilon(\eta C_G+1) + \alpha\gamma \\
&\leq& -\alpha\frac{\gamma}{2} + \alpha\gamma=\alpha\frac{\gamma}{2}<0.
\end{eqnarray*}
Concluyendo que $v\in K$. Por lo tanto 
$$q=G(\bx)+G'(\bx)(\tilde{d}+\alpha d)-v,\,v\in K,$$
lo que dice que vale la condición de Robinson.
\end{dem}

Para este caso particular, la condición de Robinson es conocida como condición de regularidad de Mangasarian-Fromovitz (MFCQ).


Ahora una condición de regularidad (más fuerte que la de Robinson y en consecuencia, más fuerte que MFCQ) muy útil para los espacios euclídeos de dimensión finita es la siguiente.

\begin{defi}
 La condición de regularidad de independencia lineal para el problema ($\tilde{P}$) (LICQ) es:
 
 $G_1'(\bx),\ldots,G_{m_\cE}'(\bx),{G_j}_1'(\bx),\ldots,{G_j}_s'(\bx)$ son linealmente independientes donde
 
 $\cA(\bx)=\{j_1,\ldots,j_s\}$ y $\bx$ en un mínimo local de ($\tilde{P}$).
\end{defi}

Esta condición de regularidad será usada en el siguiente capítulo, en la discusión sobre la convergencia del método numérico que utilizaremos para resolver el problema \eqref{pdisfinal}.

Ahora, ya teniendo los conceptos de condiciones de regularidad en espacios euclídeos, podemos dar el teorema de los multiplicadores de Lagrange para estos espacios, de manera sea posible aplicarlo a nuestro problema discretizado \eqref{pdisfinal}.

\begin{prop}[Multiplicadores de Lagrange en $\Rn$]
 Sea $\bx$ un mínimo local de f en $\Rn$, y asuma (MFCQ) vale en $\bx$. Entonces existe $\bmu\in\R^m$ tal que
\begin{equation}\label{KKTfin}
    \left\{\begin{array}{rcl}
 0&=&f'(\bx)+\sum_{i\in\cE}\bmu_i G_i'(\bx) + \sum_{j\in\cI}\bmu_j G_j'(\bx),\\
 \bmu_j\geq0 &\mbox{si}& j\in \cA(\bx),\\
 \bmu_j =0 &\mbox{si}& j\in\cI\setminus\cA(\bx).
 \end{array}\right.
\end{equation}

\end{prop}

\begin{dem}
Este teorema puede demostrarse como un corolario del teorema \eqref{lagrinf}. Por lo tanto existe $\mu\in Z^*$ tal que $$\mu\in K^* ,\quad \langle\mu , G(\bx)\rangle = 0,$$
entonces, como $K^*=\{v\in\R^m\mid v_i\in\R,\, i\in\cE;\,v_j\geq0,\,j\in\cI\}$, vale que $\mu_j\geq0$ para $j\in\cI$. Luego como $\langle\mu , G(\bx)\rangle = 0$ es fácil ver que se obtiene justamente que $\mu_j\geq0 \mbox{ si } j\in \cA(\bx)$ y $\mu_j =0 \mbox{ si } j\in\cI\setminus \cA(\bx)$.

Finalmente, por \eqref{condOptim} también vale que $$ -(f'(\bx) + G'(\bx)^*\mu)\in \cN_C(\bx).$$
Ahora por la proposición \eqref{conoconv} tenemos que $\cN_{C}(\bx)=\{v\in X^*\mid\langle v,x-\bx\rangle\leq 0\; \forall\,x\in C\}$, luego como $C=\Rn$ y $\Rn$ es un espacio vectorial, obtenemos que $\cN_{\Rn}(\bx)=\{v\in {\Rn}\mid\langle v,x\rangle= 0\; \forall\,x\in \Rn\}$ lo que implica que $-(f'(\bx) + G'(\bx)^*\mu)=0$, lo cual es lo mismo que $0=f'(\bx)+\sum_{i\in\cE}\mu_i G_i'(\bx) + \sum_{j\in\cI}\mu_j G_j'(\bx)$.
\end{dem}

La condición 
\begin{equation}\label{complementariedad}
\begin{array}{rcl}
\bmu_j\geq0 &\mbox{si}& j\in \cA(\bx),\\
 \bmu_j =0 &\mbox{si}& j\in\cI\setminus\cA(\bx).
\end{array}
\end{equation}
se conoce como condición de \texttt{complementariedad}, y se dice que se cumple la \texttt{complementariedad estricta} si vale \eqref{complementariedad} con $$\bmu_j>0 \mbox{ si } j\in \cA(\bx).$$


%Como antes mencionamos, LICQ es una condición más fuerte que la condición de Robinson, %es decir, que si vale LICQ entonces también vale la condición de Robinson. Pero tal vez %pueda ocurrir que no vale LICQ y sin embargo vale la condición de Robinson en algún %punto dado, por lo tanto el siguiente resultado nos muestra una forma más sencilla de %verificar la condición de Robinson, siempre que nustras funciones se muevan en espacios %euclídeos.



Estamos preparados entonces, para resolver el problema. Existe una gran variedad de diferentes métodos computacionales para resolver problemas de optimización, en este caso usaremos el método de \texttt{progrmación cuadrática secuencial} (SQP), por lo que necesitaremos conocer más de éste método para continuar con la resolución del problema, lo que veremos en el siguiente capítulo.

\chapter{Preliminares de métodos numéricos}

Para resolver el problema \eqref{pdisfinal} utilizaremos el método SQP de Matlab. A continuación veremos cómo funciona.
\section{Progrmación cuadrática secuencial o SQP}
Este método nos permite simular un método de Newton para optimización con restricciones, igual que como se realiza para la optimización sin restricciones \cite[Section 1.4.]{bertsekas1999nonlinear}. En cada iteración se hace una aproximación de la matriz hessiana de la función lagrangiana actualizando con un método cuasi-Newton (ver \cite[Section 1.7.]{bertsekas1999nonlinear}). Esto se usa entonces para generar un subproblema de programación cuadrática o QP (luego veremos en detalle el significado de QP) cuya solución es usada para crear una dirección de búsqueda para una línea de procedimiento de búsqueda. Ahora describiremos el método general.

Consideremos el problema $(\tilde{P})$ definido en la sección \textbf{5.2}. La idea principal es formular un subproblema de programación cuadrática QP (a continuación lo veremos en detalle) basado en una aproximación cuadrática de la función lagrangiana, ésta es 
\begin{equation}\label{funlag}
L(x,\mu)=f(x)+\sum_{i=1}^m\mu_iG_i(x).
\end{equation}

Luego se obtiene el subproblema QP linealizando las restricciones no lineales.

%\textbf{Subproblema de programación cuadrática (QP)}
%\begin{equation}\label{QP}
%\begin{array}{ll}
%    \displaystyle \mathop{\mathrm{minimizar}}_{d\in\Rn}\frac{1}{2}\la H_kd,d\ra +\la %f'(x_k),d\ra\\
%    \la g_i'(x_k),d\ra+g_i(x_k)=0,\, i=1,\ldots m_E\\
%    \la g_i'(x_k),d\ra+g_i(x_k)\leq0,\, i=m_E+1,\ldots,m
%\end{array}
%\end{equation}

\textbf{Subproblema de programación cuadrática (QP)}
\begin{equation}\label{QP}
\begin{array}{ll}
    \displaystyle \mathop{\mathrm{minimizar}}_{d\in\Rn}\frac{1}{2}d^tH_kd+f'(x_k)^td\\
    (G_i')^t(x_k)d+G_i(x_k)=0,\, i=1,\ldots m_E\\
    (G_i')^t(x_k)d+G_i(x_k)\leq0,\, i=m_E+1,\ldots,m
\end{array}
\end{equation}
Este subproblema se puede resolver usando cualquier algoritmo para QP's.

La solución es utilizada para obtener iterativamente $$x_{k+1}=x_k+\alpha_kd_k.$$
El parámetro paso $\alpha_k$ es determinado por un apropiado procedimiento de búsqueda lineal de manera que se obtenga un descenso suficiente el la función objetivo (también veremos esto a continuación). La matriz $H_k$ es una aproximación definida positiva de la matriz hessiana de la función lagrangiana \eqref{funlag}. $H_k$ se puede actualizar con cualquier método cuasi-Newton, aunque el método BFGS (lo veremos ahora) es el más popular.

La implementación del método SQP consta de tres etapas principales, las cuales discutiremos brevemente en las siguientes subsecciones:
\begin{itemize}
    \item Actualización la matriz Hessiana.
    \item Solución del problema QP.
    \item Búsqueda lineal y función objetivo.
\end{itemize}

\subsection{Actualización la matriz Hessiana}

\textbf{REVISAR ESTO}

En cada iteración definimos una aproximación $H$ cuasi-Newton de la matriz hessiana de la función lagrangiana, que se calcula mediante el método BFGS, donde $\mu_i,\,i=1,\ldots m$ es una estimación de los multiplicadores de Lagrange. Entonces

\begin{equation}\label{actmatriz}
    H_{k+1}=H_k+\frac{q_kq_k^t}{q_k^ts_k}-\frac{H_ks_ks_k^tH_k^t}{s_k^tH_ks_k},
\end{equation}
donde 
\begin{equation}\label{coefbfgs}
    \begin{array}{c}
    s_k=x_{k+1}-x_k,\\
    q_k=\left(f'(x_{k+1})+\mu^tG'(x_{k+1})\right)-\left( f'(x_k)+\mu^tG'(x_k)\right).
    \end{array}
\end{equation}

Powell (ver \cite{powell1978fast}) recomienda matener la matriz hessiana definita positiva, incluso si se vuelve indefinida positiva en el punto solución. Para hacer esto, se da $q_k^ts_k$ positivo en cada actualización y $H$ debe ser en el inicio una matriz definida positiva, entonces podemos garantizar que $H{k+1}$ será definida positiva (una demostración de esta afirmación se puede encontrar en \cite[Theorem 4.5.]{bonnans2006numerical}). La correción de Powell que realizaremos (en una versión más actualizada que la mencionada anteriormente) es la siguiente.

Modificaremos $q_k$ cuando $q_k^ts_k$ no sea lo suficientemente positivo. En una primer etapa se calula $\alpha_k$ y $d_k$ de manera que la función objetivo decrezca, lo que genera la siguiente iteración $x_{k+1}=x_k+\alpha_kd_k$. Entonces se define $q_k$ y $s_k$ por \eqref{coefbfgs} y se toma $q_k^P\in\Rn$ como una combinación convexa de $q_k$ y $H_ks_k$, esto es,
\begin{equation}\label{combpowell}
    q_k^p:=\theta q_k+(1-\theta)H_ks_k.
\end{equation}

La elección de $h_ks_k$ como vector de emergencia obtenido cuando $\theta=0$ en \eqref{combpowell} viene del hecho de que $s_k^tH_ks_k>0$ y $H_{k+1}=H_k$ si $q_k=H_ks_k$ (ver fórmula \eqref{actmatriz}). Para modificar $q_k$ lo menos posible, y preservar la mayor información posible de los datos del problema, el parámetro $\theta$ se toma tan grande como sea posible en $[0,1]$ mientras que satisfaga $$(q_k^P)^ts_k\geq\kappa s_k^tH_ks_k,$$ donde la constante $\kappa\in (0,1)$ se sugiere que sea fijada en $0.2$. Como $H_k$ se asume definida positiva, la desigualdad se satisface para $\theta=0$. Luego, un simple cálculo nos da
\begin{equation*}
    \theta=\left\{
    \begin{array}{rcl}
    1&\mbox{si}&q_k^ts_k\geq\kappa s_k^tH_ks_k,\\
    (1-\kappa)\frac{s_k^tH_ks_k}{s_k^tH_ks_k-q_k^ts_k}&\mbox{c.c}.
    \end{array}
    \right.
\end{equation*}


\subsection{Solución del problema QP}

En cada iteración del método el SQP, se resuelve un problema QP como el siguente.
\begin{equation}\label{probqp}
    \begin{array}{rcl}
    \displaystyle \mathop{\mathrm{minimizar}}_{d\in\Rn}q(d)&=&\frac{1}{2}d^tHd+c^td,\\
    A_id&=&b_i,\,i\in\cE\\
    A_id&\leq&b_i,\,i\in\cI.
    \end{array}
\end{equation}
donde $A_i$ se refiere a la i-ésima fila de la matriz $A$ de dimensión $m\times n$.

El proceso de solución envuelve dos etapas. La primera involucra el cálculo de un punto factible (si existe alguno). La segunda etapa trata la creación de una sucesión de puntos factibles que converjan a la solución. En este método, un conjunto activo $\bar A_k$, es una estimación de las restricciones activas (es decir aquellas que se encuentran en la frontera) en el punto solución.

$\bar A_k$ se actualiza en cada iteración $k$ y se usa para formar una base para una direccion de búsqueda $\hat{d}_k$. Las restricciones de igualdad siempre se mantienen en el conjunto $\bar A_k$. La notación de la variable $\hat{d}_k$ es usada para distinguirla de la $d_k$ en la iteraciones del método SQP. La direccion de búsqueda $\hat{d}_k$ es calculada y minimiza la función objetivo siempre que se encuentre en la frontera de alguna restricción activa. El subespacio factible para $\hat{d}_k$ es formado a partir de una base $Z_k$ cuyas columnas son otrogonales a la estimación del conjunto activo $\bar A_k$ (es decir $\bar A_kZ_k=0$). Por lo que una dirección de búsqueda, la cual es formada a través de una suma lineal de cualquier combinación de las columnas de $Z_k$, tienen garantizado permanecer en la frontera de las restricciones activas.

La matriz $Z_k$ se forma desde las últimas $m-l$ columnas de la descomposición QR de la matriz $\bar A_k^t$, donde $l$ es el número de restricciones activas y $l<m$. Entonces, $Z_k$ está dado por $$Z_k=Q[:,l+1:m],$$
donde $$Q^t\bar A_k^t=[R;0].$$

Una vez que obtenemos $Z_k$, se intenta encontrar una dirección de búsqueda que minimice  $q(d)$ donde $\hat{d}_k$ está en el espacio nulo de las restricciones activas. Esto es, $\hat{d}_k$ es una combinación lineal de las columnas de $Z_k$, $\hat{d}_k=Z_kp$ para algún vector $p$.
Entonces si vemos $q$ como una función de $p$, al sustituir por $\hat{d}_k$ obtenemos $$q(p)=\frac{1}{2}p^tZ_k^tHZ_kp+c^tZ_kp.$$
Derivando con respecto a p, se tiene $$q'(p)=Z_k^tHZ_kp+Z_k^tc.$$
$\nabla q(p)$ es el gradiente proyectado de la función cuadrática, porque es el gradiente proyctado en el espacio definido por $Z_k$. El término $Z_k^tHZ_k$ se denomina hessiano proyectado. Asumiendo que la matriz hessiana $H$ es definida positiva (este es el caso en la implementación del método SQP), entonces el mínimo de la función $q(p)$ en el subespacio definido por $Z_k$ se encuentra cuando $\nabla q(p)=0$, la cual es la solución del sistema de ecuaciones lineales $$Z_k^tHZ_kp=-Z_k^tc.$$
Luego, un paso se da de la forma $$x_{k+1}=x_k+\alpha\hat{d}_k,$$
donde $\hat{d}_k=Z_kp$.

Por la naturaleza cuadrática de la función objetivo, en cada iteración, hay sólo dos opciones para el tamaño del paso $\alpha$. Un paso en una unidad en $\hat{d}_k$ es el paso exacto al mínimo de la función restringida al espacio nulo de $\bar A_k$. Si tal paso se puede dar, sin violar las restricciones, entonces ésta es la solución del problema QP \eqref{probqp}. En el otro caso, el paso en $\hat{d}_k$ a la restricción más cercana es menor que una unidad y una nueva restricción se incluye en el conjunto activo en la próxima interación. La distancia a la frontera de las restricciones está dada por $$\alpha=\displaystyle \mathop{\mathrm{minimizar}}_{i\in\{1,\ldots,m\}}\frac{-(A_ix_k-b_i)}{A_i\hat{d}_k},$$
el cual está definido por las restricciones que no están en el conjunto activo, y donde la dirección $\hat{d}_k$ apunta hacia la frontera de la restricción, es decir, $A_i\hat{d}_k>0,\,i=1,\ldots m$.

Cuando $n$ restricciones independientes se incluyen en el conjunto activo, sin la localización del mínimo, los multiplicadores de Lagrange, $\mu_k$ se calculan de manera que satisfagan el conjunto de ecuaciones lineales no singulares $$\bar A_k^t\mu_k=c.$$
Si todos los elementos de $\mu_k$ son positivos, $x_k$ es la solución optimal del problema \eqref{probqp}. Sin embargo, sin cualquier componente de $\mu_k$ es negativo, y el componente no corresponde a una restricción de igualdad, entonces el correspondiente elemento es eliminado del conjunto activo y una nueva iteración es buscada.

\textbf{Inicialización}

El algoritmo requiere un punto factible para comenzar. Si el punto actual del método SQP no es factible, se puede encontrar un punto que si lo sea resolviendo el problema de programación lineal
\begin{equation}\label{eqgamma}
    \begin{array}{ll}
   \displaystyle \mathop{\mathrm{minimizar}}_{y\in\R,\,x\in\Rn}\,\,\gamma\mbox{ tal que}\\
    A_ix=b_i,\,i\in\cE\\
    A_ix\leq b_i,\,i\in\cI.
    \end{array}
\end{equation}
Como antes, la notación $A_i$ indica la i-ésima fila de la matriz $A$. Se puede encontrar un punto factible (si existe) a la ecuación \eqref{eqgamma} poniéndole un valor a $x$ que satisfaga las restricciones de igualdad. Entonces se puede determinar este valor resolviendo un sistema de ecuaciones lineales formado del conjunto de restricciones de igualdad. Si existe una solución a este problema, la variable de holgura $\gamma$ se establece en la máxima restricción de igualdad en este punto.


\textbf{QUE ES UN PROBLEMA LP}

Se puede modificar el algoritmo QP anterior por problemas LP haciendo la dirección de búsqueda una dirección de descenso cada vez más empinada en cada iteración, aquí $\varrho_k$ es el gradiente de la función objetivo. $$\hat{d}_k=-Z_kZ_k^t\varrho_k.$$
Si un punto factible se encuentra usando este método LP, se ingresa a la etapa QP principal. La dirección de búsqueda $\hat{d}_k$ se inicia con una dirección $\hat{d}_1$ encontrada resolviendo el sistema de ecuaciones lineales $$H\hat{d}_1=-\varrho_k,$$
donde $\varrho_k$ es el gradiente de la función objetivo en la actual iteración $x_k$ (es decir, $H_{x_k}+c$).
Si una solución factible no se encuentra para el problema QP, la dirección de búsqueda para la rutina SQP principal $\hat{d}_k$ se toma como la que minimiza $\gamma$.

\subsection{Búsqueda lineal y función objetivo}

La solución del subproblema QP produce un vector $d_k$ el cual es usado para formar una nueva iteración $$x_{k+1}=x_k+\alpha_kd_k.$$
El tamaño del paso (parámetro $\alpha_k$) se determina en orden de producir un decrecimiento suficiente en la función objetivo. La función objetivo de la siguiente forma se usa en esta implementación. $$\psi(x)=f(x)+\sum_{i\in\cE}r_i\varrho_i(x)+\sum_{i=1}^mr_i\mathop{\mathrm{max}}[0,\varrho_i(x)].$$
Powel recomienda ajustar el parámetro de penalidad $$r_i=(r_{k+1})_i=\mathop{\mathrm{max}}_i\left\{\mu_i,\frac{(r_k)_i+\mu_i}{2}\right\},\, i=1,\ldots,m.$$
Esto permite la contribución de las restricciones que están inactivas en la solución del problema QP pero que estuvieron recientemente activas. En esta implemetación, el parámetro de penalidad $r_i$ se elige inicialmente como $$r_i=\frac{\|f'(x)\|_{\Rn}}{\|\varrho'(x)\|_{\Rn}}.$$
Esto asegura mayores contribuciones del parámetro de penalidad para restricciones con gradientes más pequeños, el cual sería el caso para restricciones activas en el punto solución.

\section{Algoritmo}

\textbf{Notación}. Dado un vector $v\in\R^m$, lo asociamos al vector $v^@\in\R^m$, que se define como sigue.
\[(v^@)_i=
\left\{\begin{array}{rcl}
v_i &\mbox{si}& i\in \cE\\
v_i^+&\mbox{si}& i\in \cI,
\end{array}\right.
\]
donde $v_i^+=\mathrm{max}\{0,v_i\}$.


Una manera de implementar el método SQP es seguir los siguientes pasos.
 
 \begin{itemize}
     \item Se da un punto inicial para iterar $(x_1,\mu_1)$.
     \item Calculamos $G(x_1),G'(x_1)$ y $f'(x_1)$.
     \item Comenzamos con $k=1$.
 \end{itemize}
 \begin{enumerate}
 \item Parar si se cumplen las condiciones KKT \eqref{KKTfin} en $(\bx,\bmu)=(x_k,\mu_k)$ (es decir, la optimalidad es alcanzada).
 \item Calcular $H_k$ y encontrar un punto estacionario para el problema \eqref{QP} dicha solución es la misma que la del siguiente sistema modificado.
 \begin{equation}\label{qpmodificado}
     \begin{array}{ll}
 H_kd+(G'(x_k))^t\mu^{QP}=-f'(x_k),\\
 (G(x_k)+(G'(x_k))^td)^@=0,\\
 (\mu^{QP})_\cI\geq0,\\
  (\mu^{QP})_\cI^t(G(x_k)+(G'(x_k))^td)_\cI=0,
 \end{array}
 \end{equation}
 a la cual denotaremos $(d_k,\mu_k^{QP})$
 \item Definir $x_{k+1}=x_k+d_k$ y $\mu_{k+1}=\mu^{QP}$.
 \item Calcular $G(x_{k+1}),f'(x_{k+1})$ y $H_{k+1}$.
 \item Ajustar $k$ a $k+1$ y volver al paso 1.
 \end{enumerate}
 
 Con la información de este capítulo, ya somos capaces de implementar el método SQP conociendo en detalle su funcionamiento. En la siguiente sección, veremos algunos resultados sobre la convergencia de este método, lo que nos mustra una de las razones por la cual el SQP funciona, y hemos optado por ultilizarlo.
 
 \section{Convergencia del método}

El teorema de convergencia que veremos a continuación resalta una interesante propiedad del algoritmo SQP, ésta es, que en algún entorno de un punto estacionario regular (veremos esta noción en detalle en el teorema) que satisface estrictamente la complementariedad, las restricciones activas del problema QP \eqref{QP} son las mismas que las del problema $(\tilde{P})$. El resultado que veremos más abajo dice que las restricciones activas de un problema QP en una iteración, probablemente sean las mismas en la próxima iteración, al menos cercanas a un punto estacionario regular. Numéricamente, esto significa que es de utilidad resolver los problemas QP con algoritmos que tomen ventaja de una buena suposición de las restricciones activas. Entonces, el problema combinatorio de determinar cuáles son las restricciones activas en la solución ya no ocurren durante las últimas iteraciones del algoritmo SQP.

\textbf{Complementariedad estricta y notacion en el siguiente teorema y superlinealidad y cuadratico?  Citar el otro libro de Bonnans QUE ES UN PUNTO DE NORMA MÍNIMA}


Para el teorema de convergencia del método SQP necesitaremos el siguiente resultado.

\begin{teo}[Convergencia del algoritmo de Newton]\label{convernewton}
Consideremos el problema $\tilde{P}$ sólamente con restricciones de igualdad (lo denominaremos $\tilde{P}_E$. Suponga que $f$ y $G$ son de clase $C^2$ en un entorno de un punto regular estacionario $\bx$ de el problema $\tilde{P}_E$, con multiplicador asociado $\bmu$. Entonces, existe un entorno $V$ de $(\bx,\bmu)$ tal que, si la primer iteración $(x_1,\mu_1)\in V$, el algoritmo de Newton (ver \cite[Section 14.1]{bonnans2006numerical}) está bien definido y genera una sucesión $\{(x_k,\mu_k)\}$ que converge superlinealmente a $(\bx,\bmu)$.
\end{teo}

La demostración del teorema anterior se puede encontrar en \cite[Therorem 14.4]{bonnans2006numerical}.

\begin{teo}
Suponga que $f$ y $g$ son de clase $C^2$ en un entorno de un punto estacionario $\bx$ del problema $(\tilde{P})$, con un multiplicador asociado $\bmu$. Supunga también que se cumple la complementariedad estricta y que $(\bx,\bmu)_{E\cup \cA(\bx)}$ es un punto estacionario regular del problema con restricciones de igualdad 
\begin{equation}\label{pconver}
    \begin{array}{ll}
    \displaystyle \mathop{\mathrm{minimizar}}_{x} f(x)\\
    G_i(x)=0, \mbox{    para } i\in \cE\cup \cA(\bx)
    \end{array}
\end{equation}
donde un punto estacionario $(\bx,(\bmu)$ de $(\tilde{P})$ se dice regular, si $A_*:=g'(\bx)$ es sobreyectiva y $Z_*^tL_*Z_*^-$ es no singular, para alguna (o algunas) matriz $Z_*^-$ $n\times (n-m)$ cuyas columnas forman una base de $\mathrm{Ker}(A_*)$.

Considere el algoritmo SQP, en el que $d_k$ es un punto estacionario de mínima norma del problema QP \eqref{QP}. Entonces existe un entorno $V$ de $(\bx,(\bmu)$ tal que, si la primer iteración $(x_1,\mu_1)\in V$,
\begin{enumerate}[i)]
\item el algoritmo está bien definido y genera una sucesión $\{(x_k,\mu_k)\}$ que converge superlinealmente a $(\bx,\bmu)$,
\item las restricciones activas del problema QP \eqref{QP} son las del problema $(\tilde{P})$,
\end{enumerate}
\end{teo}
\begin{dem}
La idea de la prueba es mostrar que, cerca de $(\bx,\bmu)$ el punto estacionario de mínima norma del problema QP \eqref{QP} y el paso de Newton para \eqref{pconver} son idénticos. Entonces el resultado sigue del teorema \eqref{convernewton}.

Supongamos que $(x,\mu)$ está cerca de $(\bx,\bmu)$. Como $(\bx,(\bmu_{\cI\cup\cA(\bx)})$ es un punto estacionario regular de \eqref{pconver}, $G'_{\cI\cup\cA(\bx)}(\bx)$ es sobreyectiva y el problema en $\tilde{d}$
\begin{equation}\label{qpdemteor}
    \begin{array}{l}
    \mathrm{min}_{\tilde{d}}\ f'(x)^t\tilde{d}+\frac{1}{2}\tilde{d}L(x,\mu)\tilde{d},\\
    G_i(x)+G_i'(x)^t\tilde{d}=0,\quad \mbox{para }i\in\cI\cup\cA(\bx),
    \end{array}
\end{equation}
tiene un único punto estacionario, al cual lo denotamos por $(\tilde{d},\tilde{\mu}_{\cI\cup\cA(\bx)})$, y formamos con $\tilde{\mu}_{\cI\cup\cA(\bx)}$ un vector $\tilde{\mu}\in\R^m$, haciendo $\tilde{\mu}_i=0$ para $i\in\cI\setminus\cA(\bx)$.

Ahora, mostraremos que $(\tilde{d},\tilde{\mu})$ es un punto estacionario del problema QP \eqref{QP}, si $(x,\mu):=(x_k,\mu_k)$ está en algún entorno de $(\bx,\bmu)$. Sólo necesitamos ver que para $G_i(x)+G_i'(x)^t\tilde{d}\leq0$ para $i\in\cI\setminus\cA(\bx)$ y $\mu_i\geq0$ para $i\in\cA(\bx)$. Por el teorema \eqref{convernewton}, $(x+\tilde{d},\bmu)$ está cerca de $(\bx,\bmu)$ cuando $(x,\mu)$ está cerca de $(\bx,\bmu)$. Por lo tanto, para $i\in\cA(\bx)$, $\tilde{\mu}_i\geq0$, a que $(\bmu)_i>0$ por la complemetariedad estricta. Por otro lado $\tilde{d}$ es pequeño, de manera que $G_i(x)+G_i'(x)^t\tilde{d}\leq0$ para $i\in\cI\setminus\cA(\bx)$. Por lo tanto $(\tilde{d},\tilde{\mu})$ es un punto estacionario de \eqref{QP}. De esto deducimos que, para $(x,\mu)$ cerca de $(\bx,\bmu)$ el algoritmo SQP está bien definido, $d$ es un punto estacionario de norma mínima y $\tilde{d}$ es pequeño por el el teorema \eqref{convernewton}.

Veamos ahora que el par $(d,\mu^{QP}):=(d_k,\mu^{QP}_k)$ formado por la solución de mínima norma del problema QP y su multiplicador asociado es, de hecho, $(\tilde{d},\tilde{\mu})$, si $(x,\mu)$ está en algpun entorno de $(\bx,\bmu)$. Por el teorema \eqref{convernewton} concluiría la demostración. 

Para $(x,\mu)$ cerca de $(\bx,\bmu)$ y $i\in\cI\setminus\cA(\bx)$, $G_i(x)+G_i'(x)^td<0$, entonces $\mu_i^{QP}=0=\tilde{\mu}_i$. Por la unicidad del punto estacionario de \eqref{qpdemteor}, basta mostrar que  $G_i(x)+G_i'(x)^td=0$ para $i\in\cA(\bx)$ y $(x,\mu)$ cerca de  $(\bx,\bmu)$. Si este no es el caso, debe existir un índice $j\in\cA(\bx)$ y una sucesión $(x,\mu)\to(\bx,\bmu)$ (\textbf{QUE SIGNIFICA ESTO}), tal que $G_j(x)+G'_j(x)^td<0$. Entonces $\mu_j^{QP}=0$ y $$f'(x)+L(x,\mu)d+\sum_{i\in\cI\cup\cA(\bx)\setminus \{j\}} \ \mu_i^{QP}G_i'(x)=0.$$
Como $f'(x)+L(x,\mu)d\to f'(\bx)$ y $G'_{\cI\cup\cA(\bx)}(\bx)$ es sobreyectiva, $\mu_i^{QP}$ para $i\in\cI\cup\cA(\bx)\setminus \{j\}$ converge a algún límite, el cual llamaremos $\mu_i^*$. Tomando límite en la ecuación de arrtiba obtenemos $$f'(\bx)+\sum_{i\in\cI\cup\cA(\bx)\setminus \{j\}} \ \mu_i^*G_i'(\bx)=0.$$
Por lo tanto encontramos dos multiplicadores diferentes: $\mu^*$ (hacemos $\mu_i^*=0$ para $i\in\cI\cup\cA(\bx)\setminus \{j\})$ y $\bmu$ (son distintos pues $\mu^*_j=0$ y $\bmu_j>0$ por la complemetariedad estricta). Esto genera una contradicción con la unicidad de los multiplicadores, que vale por la sobreyectividad de $G'_{\cI\cup\cA(\bx)}(\bx)$.
\end{dem}

El teorema anterior vale si no exigimos complementariedad estricta y en cambio se asume que $(\bx,\bmu)$ satisfacen condiciones de optimalidad de segundo orden e independencia lineal de los gradientes de las restricciones activas (LICQ). Este resultado se puede ver en detalle en \cite[Theorem 15.4]{bonnans2006numerical}.

Con esto finalizamos el Capítulo y comenzamos a continuación a analizar los resultados numéricos obtenidos utilizando el método SQP de MATLAB.


%%%%%%%%%%%%%%%%%%%MIRAME

\chapter{Planteo final del problema}

\section{Dimensión infinita}
\subsection{Notación de utilidad}

Para ser consistentes con la notación del teorema de los multiplicadores de Lagrange, reescribimos el problema \eqref{problema} como
\begin{equation}\label{pfinal}
\begin{array}{rl}
\displaystyle \mathop{\mathrm{minimizar}}_{x} & \displaystyle f(x) \\
\mathrm{sujeto \ a} & G(x)\in K,\, x\in C
\end{array}
\end{equation}

donde $f:X\arr \R$ con 
\[
 X=H^2[0,1]^2\times L^2[0,1]\times \R_+ \times \R_+,
\]
tal que para $x=(c,a,p,v_0)\in X$,
$$
 f(x)=\int_{0}^{1} \frac{p}{v(s)}ds.
$$
Las restricciones $G:X\mapsto Z$ donde
\[
Z=L^2[0,1]\times \R^2\times \R\times \R\times \R^2\times \R^2\times L^2[0,1] \times L^2[0,1]^2 \times \R,
\]
y por lo tanto $G=(G_1,G_2,G_3,G_4,G_5,G_6,G_7,G_8,G_9)$, tal que
\begin{eqnarray*}
 G_1(x)(s)&=&\left\|\frac{dc}{ds}(s)\right\|_{\R^2}^2-1,\quad \mbox{p.p. }s\in(0,1)\\
 G_2(x)&=&c(1)-c(0),\\
 G_3(x)&=&\int_{0}^{1} a(r)dr,\\
 G_4(x)&=&a(1)-a(0),\\
 G_5(x)&=&\frac{dc}{ds}(1)-\frac{dc}{ds}(0),\\
 G_6(x)&=&\frac{d^2c}{ds^2}(1)-\frac{d^2c}{ds^2}(0),\\
 G_7(x)(s)&=&\frac{1}{\mu_t^2} a(s)^2 + \frac{1}{\mu_n^2 p^2}v(s)^4\left\| \frac{d^2 c}{ds^2}(s) \right\|_{\R^2}^2-g^2,\quad \mbox{p.p. }s\in(0,1)\\
 G_8(x)(s)&=&\|pc(s)-\alpha(s)\|_{\R^2}^2-\kappa(s)^2, \quad \mbox{p.p. }s\in(0,1)\\
 G_9(x)(s)&=&\|pc(s)-\beta(s)\|_{\R^2}^2-\kappa(s)^2, \quad \mbox{p.p. }s\in(0,1).\\
\end{eqnarray*}
Adem\'as, $C\subset X$ y $K\subset Z$, dados por
\begin{eqnarray*}
C&=&H^2[0,1]^2\times L^2[0,1]\times \R_+ \times \R_+,\\
K&=&\{0\}\times \{0\}\times \{0\}\times \{0\}\times \{0\}\times \{0\}\times (-L^2[0,1]_+) \times (-L^2[0,1]_+) \times (-L^2[0,1]_+).
\end{eqnarray*}


\subsection{Cálculo de derivadas}
En esta parte del trabajo realizaremos el cálculo de las derivadas de las funciones que definen el problema \eqref{problema}. Para ello, notemos que si $F:X\arr Y$ es Fréchet diferenciable, entonces $F$ es direccionalmente diferenciable y podemos calcular $F'(\bx)$ usando que para toda dirección $\tilde x\in X$,
$$F'(\bx)(\tilde x)=\left.\frac{d}{dt}F(x+t\tilde{x})\right|_{t=0}.$$

Por lo tanto, sea $\tilde{x}=(\tilde{c},\tilde{a},\tilde{p},\tilde{v_0})$ una dirección arbitraria en $X$. Entonces para la función objetivo $f:X\arr\R$,
\begin{eqnarray*}
 \langle f'(x), \tilde{x}\rangle
&=&\frac{d}{dt}\left.\left( \int_{0}^{1}\textstyle \frac{p+t\tilde{p}}{\sqrt{2(p+t\tilde{p}) \int_0^s (a(r)+t\tilde{a}(r))dr + (v_0+t\tilde{v_0})^2}}ds\right)\right|_{t=0}\\
&=&\int_0^1\left(-\frac{p^2}{v(s)^3}\right)\int_0^s\tilde{a}(r)dr ds + \int_0^1\left(\frac{1}{v(s)}-\frac{p}{v(s)^3}\int_0^s a(r)dr\right)ds\, \tilde{p} + \int_0^1\left(-\frac{p v_0}{v(s)^3}\right)ds\, \tilde{v}_0.
\end{eqnarray*}
Para la restricción $G:X\arr Z$ las derivadas de las funciones coordenadas son
\begin{eqnarray*}
G_1'(x)(\tilde{x})
&=&2\left\langle\frac{dc}{ds}(s),\frac{d\tilde{c}}{ds}(s)\right\rangle,\\
G_2'(x)(\tilde{x})&=&\tilde{c}(0)-\tilde{c}(1),\\
G_3'(x)(\tilde{x})&=&\int_0^1 \tilde{a}(r)ds,\\
G_4'(x)(\tilde{x})&=&\tilde{a}(0)-\tilde{a}(1),\\
G_5'(x)(\tilde{x})&=&\frac{d\tilde{c}}{ds}(0)-\frac{d\tilde{c}}{ds}(1),\\
G_6'(x)(\tilde{x})&=&\frac{d^2\tilde{c}}{ds^22}(0)-\frac{d^2\tilde{c}}{ds^2}(1),\\
G_7'(x)(\tilde{x})&=&\frac{2a(s)\tilde{a}(s)}{\mu_t^2} + \frac{1}{\mu_n^2}\left[-\frac{2p\tilde{p}}{p^4}v(s)^4\left\|\frac{d^2c}{ds^2}(s)\right\|_{\R^2}^2\right.\\
&&+\frac{1}{p^2}\left(2v(s)^2(2\tilde{p}\int_0^s a(r)dr +2p\int_0^s\tilde{a}(r)dr+2v_0\tilde{v}_0) \left\|\frac{d^2c}{ds^2}(s)\right\|_{\R^2}^2\right.\\
&&\left.\left.+ 2v(s)^4\left\langle\frac{d^2c}{ds^2}(s) , \frac{d^2\tilde{c}}{ds^2}(s)\right\rangle_{\R^2}\right)\right],\\
G_8'(x)(\tilde{x})&=&\langle2p(pc(s)-\alpha(s)) , \tilde{c}(s)\rangle_{\R^2} + 2\tilde{p}\langle(pc(s)-\alpha(s)) , c(s)\rangle_{\R^2},\\
G_9'(x)(\tilde{x})&=&\langle2p(pc(s)-\beta(s)) , \tilde{c}(s)\rangle_{\R^2} + 2\tilde{p}\langle(pc(s)-\beta(s)) , c(s)\rangle_{\R^2}.\\
\end{eqnarray*}

Aunque ya tenemos el planteo final, todavía no podemos resolver el problema, ya que por sus características, es demasiado complejo. Así que necesitaremos utilizar la computadora y métodos numéricos para resolver el problema, pero no podemos trabajar de esta manera en dimensión infinita, por lo que llevaremos todo a dimensión finita mediante aproximaciones. 

\section{Dimensión finita}
\subsection{Discretización del problema}

Para plantear el problema de una forma que la computadora comprenda, queremos aproximar nuestros espacios y funciones por alguna versión de ellos discretizada. La idea entonces es usar el \texttt{método de elementos finitos}, en el cual elegimos una base de funciones finita para nuestros espacios y aproximamos sus elementos con coeficientes reales.

Aproximaremos al espacio 
$X=H^2[0,1]^2\times L^2[0,1]\times \R \times \R$\quad por $\quad\textbf{X}=V^2\times V\times \R\times \R$,
aquí $\, V=\mbox{span}\left\{\varphi_j,\ j=1,\ldots, n\right\}$.
Donde tomaremos las siguientes funciones base (ver \cite[Subsection 2.3.1]{stoer1993introduction}),
\begin{eqnarray*}\label{funbase}
\varphi_j(t)&=&\left\{ \begin{array}{ll}
\frac{1}{2}, & \mbox{ si}\quad j=1,\\
\cos(2\pi kt), & \mbox{ si}\quad j=2k+1,\\
\mathrm{sen}(2\pi kt), & \mbox{ si}\quad j=2k,
\end{array}
\right.\mbox{para } k=1,\ldots,m,
\end{eqnarray*}
con $n=2m+1$. Aproximaremos al espacio $Z$ de la siguiente manera.
Para las coordenadas de $Z$ que se encuentran en $\R$, la función base que usaremos es el número real $1$, pues span$\{1\}=\R$. En las coordenadas funcionales de $Z$ aproximaremos las funciones por su valor en una malla de $n+1$ puntos entre $0$ y $1$ dados por $s_i=\frac{i-1}{n}$ $i=1,\ldots,n+1$, para esto, utilizaremos la función $\delta_{s_i}:C_c^\infty[0,1]\arr\R$, de manera que $\la\delta_{s_i},f\ra=f(s_i)$ donde $f$ es una función a valores reales.

\textbf{Observación}. Como el espacio $C_c^\infty[0,1]$ es denso en $L_2[0,1]$, es posible usar las funciones $\delta_{s_i}$ en nuestro problema.
Por otro lado, la elección de las funciones base tienen que ver con su propiedad de periodicidad, la cual facilitará las cuentas en el planteo discretizado.

Será útil, para probar que nuestro problema discretizado tiene solución, ver que las funciones base son linealmente independientes. Para esto utilizaremos números complejos, que facilitan notación y cuentas en este caso.

Consideremos entonces $s_j$ con $j=1,\ldots,n$ (por periodicidad sabemos que las funciones base conciden en los puntos $s_1=0$ y $s_{n+1}=1$). Ahora definimos

\begin{equation}\label{wjdef}
    w_j=e^{i2\pi s_j}.
\end{equation}
En consecuencia, podemos observar que $w_j^n=e^{i2\pi(j-1)}=1$ ya que $j-1$ es entero. Luego, obtenemos que 
\begin{equation}\label{wjprop1}
    w_j^{-k}=e^{-i2\pi ks_j}=1.(e^{-i2\pi ks_j})=(e^{i2\pi ns_j})(e^{-i2\pi ks_j})=e^{i2\pi ns_j-i2\pi ks_j}=e^{i2\pi s_j(n-k)}=w_j^{n-k}.
\end{equation}


Una observación importante es que 
\begin{equation}\label{wjkprop2}
    w_j^k=e^{i2\pi k\frac{j-1}{n}}=e^{i2\pi\frac{k+1-1}{n}(j-1)}=w_{k+1}^{j-1},
\end{equation}
para $j=1,\ldots,n$ y $k=0,\ldots n-1$.

Ahora, tomando $0\leq h\leq j\leq n-1$ tenemos que $w_{j+1-h}^n=1$ y esto implica que $\mathrm{cos}(2\pi\frac{j-h}{n})=1$ por lo que $j=h$, luego usando \eqref{wjkprop2} y \eqref{wjprop1}, 

\textbf{por qué vale esto:}

\begin{equation}\label{wjproducto}
    0=\displaystyle\sum_{k=0}^{n-1}w_{j-1-h}^k=\displaystyle\sum_{k=0}^{n-1}w_{k+1}^{j-h}=\displaystyle\sum_{k=1}^{n}w_k^jw_k^{-h},
\end{equation}
por lo tanto
\begin{equation}\label{wkproducto}
    \displaystyle\sum_{k=1}^{n}w_k^jw_k^{-h}=\left\{\begin{array}{ll}
         n&\mbox{si }j=h\\
         0&\mbox{si }j\neq h.
    \end{array}\right.
\end{equation}

Entonces, usaremos estos números complejos $w_j$ dados por \eqref{wjdef} para dar una expresión de una combinación lineal de escalares con nuestras funciones base, de manera que logremos demostrar la independencia lineal utilizando \eqref{wjproducto}.

Sean $\lambda_1,\ldots,\lambda_n\in\R$, entonces podemos escribir
\begin{equation}\label{coefli}
    \begin{array}{rcl}
    \displaystyle\sum_{h=1}^n\lambda_h\varphi_h(s_j)&=&\frac{\lambda_1}{2}+\displaystyle\sum_{k=1}^m\lambda_{2k}\varphi_{2k}(s_j)+\lambda_{2k+1}\varphi_{2k+1}(s_j)\\
     &=&\frac{\lambda_1}{2}+\displaystyle\sum_{k=1}^m\lambda_{2k}\frac{w_j^k-w_j^{-k}}{2i}+\lambda_{2k+1}\frac{w_j^k+w_j^{-k}}{2}\\
     &=&\frac{\lambda_1}{2}+\displaystyle\sum_{k=1}^m\frac{\lambda_{2k+1}-i\lambda_{2k}}{2}w_j^k+\frac{\lambda_{2k+1}+i\lambda_{2k}}{2}w_j^{n-k}\\
     &=&\frac{\lambda_1}{2}+\displaystyle\sum_{k=1}^m\frac{\lambda_{2k+1}-i\lambda_{2k}}{2}w_j^k+\displaystyle\sum_{r=m+1}^{n-1}\frac{\lambda_{2(n-r)+1}+i\lambda_{2(n-r)}}{2}w_j^r\\
     &=&\displaystyle\sum_{k=0}^{n-1}\zeta_kw_j^k,
    \end{array}
\end{equation}
donde
\begin{equation*}
    \begin{array}{rcl}
    \zeta_0&=&\frac{\lambda_1}{2},\\
    \zeta_k&=&\frac{\lambda_{2k+1}-i\lambda_{2k}}{2}\mbox{ para }k=1,\ldots,m;\\
    \zeta_k&=&\frac{\lambda_{2(n-k)+1}+i\lambda_{2(n-k)}}{2}\mbox{ para }k=m+1,\ldots,n-1.
    \end{array}
\end{equation*}

Para concluir la prueba, tomemos $0\leq k,h\leq n-1$ y supongamos que (usando \eqref{coefli}) $$\displaystyle\sum_{h=1}^n\lambda_h\varphi_h(s_j)=\displaystyle\sum_{k=0}^{n-1}\zeta_kw_j^k=0,$$
entonces
\begin{equation*}
    0=\displaystyle\sum_{j=1}^n\left(\displaystyle\sum_{k=0}^{n-1}\zeta_kw_j^k\right)w_j^{-h}=\displaystyle\sum_{k=0}^{n-1}\zeta_kn\delta_{kh}=n\zeta_h=\left\{\begin{array}{lcl}
     n\frac{\lambda_1}{2}&\mbox{si}&h=0,\\
     n\frac{\lambda_{2h+1}-i\lambda_{2h}}{2}&\mbox{si}&1\leq h\leq m,\\
     n\frac{\lambda_{2(n-h)+1}+i\lambda_{2(n-h)}}{2}&\mbox{si}&m+1\leq h\leq n-1.
\end{array}\right.
\end{equation*}
donde 
\begin{equation*}
    \delta_{ij}=\left\{\begin{array}{lcl}
        1&\mbox{si}&i=j,  \\
        0&\mbox{si}&i\neq j.
    \end{array}\right.
\end{equation*}

Por lo tanto $\lambda_1=0$, y $\frac{\lambda_{2h+1}-i\lambda_{2h}}{2}=0,$ entonces $$\lambda_{2h+1}=i\lambda_{2h},$$
luego, como tanto la parte real como la imaginaria deben coincidir, se tiene que $$0=\lambda_{2h+1}=\lambda_{2h},\mbox{ para }h=1,\ldots,m;$$
esto implica que $\lambda_i=0, \mbox{ para }i=1,\ldots,n$, concluyendo que las funciones base son linealmente idependientes.




Entonces, teniendo en cuenta el planteo del problema \eqref{pfinal}, podemos comenzar a reescribir nuestras funciones. Primero nuestra variable $x$ que consta de cinco coordenadas, a la cual aproximamos en las coordenadas cuyos elementos son funciones por una suma finita de coeficientes multiplicados por nuestras funciones base, y como las demás coordenadas son números reales, no es necesario discretizarlas. Por lo tanto queda $$x=(c_1,c_2,a,p,v_0)\approx\hat{x}=\left(\sum_{k=1}^n\texttt{c}^1_k\varphi_k, \sum_{k=1}^n\texttt{c}^2_k\varphi_k,\sum_{i=1}^n\texttt{a}_i\varphi_i, p, v_0\right),$$
entonces en nuestro problema discretizado, nuestra nueva variable real es $$\texttt{x}=(\texttt{c}^1,\texttt{c}^2,\texttt{a},p,v_0)\in\R^{3n+2}.$$

Así, podemos aproximar nuestra función $v$ por una nueva función $\texttt{v}$ reemplaando de manera análoga su aceleración tangencial $a$ por su aproximación con coeficientes $\texttt{a}$, matemáticamente $$\texttt{v}(s)=\sqrt{2p\int_0^s \sum_{i=1}^n\texttt{a}_i\varphi_i(r)dr +v_0^2},$$
de la misma manera procedemos a aproximar la función objetivo $f$, obteniendo $$f(x)\approx\texttt{f}(\texttt{x}):=f(\hat x)=\int_0^1 \frac{p}{\textbf{v}(s)}ds.$$

Ahora debemos aproximar también las funciones restricciones, para ello separamos dos casos, las restricciones que llegan a los números reales y las que llegan al espacio $L^2[0,1]$. Entonces recordando las restricciones del problema \eqref{pfinal} vemos que $G_2, G_3, G_4, G_5$ y $G_6$ llegan a $\R$ y las restantes llegan a $L^2[0,1]$. Luego a las restricciones que llegan a $\R$ podemos aproximarlas de la misma manera que hicimos antes con nuestra función objetivo, por lo que tenemos $$\texttt{G}_r(\texttt{x}):=G_r(\hat x),\quad r=2,3,4,5,6.$$
Ahora para las restricciones que llegan a $L^2[0,1]$ usamos la antes mencionada función $\delta$, por lo tanto obtenemos $$(\texttt{G}_r)_{i}(\texttt{x}):=\langle G_r(\hat x) , \delta_{s_i}\rangle,\,i=1,\ldots,l; \quad r=1,7,8,9.$$

Notemos que, como vimos en la observación anterior, nuestras funciones base son periódicas con periodo $1$, es decir que $\varphi_j(0)=\varphi_j(1)\, \forall j$ (simplemente por propiedades trigonométricas), lo que implica que $\texttt{G}_2,\texttt{G}_4,\texttt{G}_5\mbox{ y }\texttt{G}_6$ son iguales a 0, por lo que esas restricciones se cumplen automáticamente.

Luego, el planteo del problema queda así

\begin{equation}\label{pdisfinal}
\begin{array}{rl}
\displaystyle \mathop{\mathrm{minimizar}}_{(\texttt{c}^1,\texttt{c}^2,\texttt{a},p,v_0)} & \texttt{f}(\texttt{x})=\int_0^1\frac{p}{\sqrt{2p\int_0^s \sum_{i=1}^n\texttt{a}_i\varphi_i(r)dr +v_0^2}}\\
 \mathrm{sujeto \ a} & \left\|\left(\sum_{k=1}^n\texttt{c}^1_k\varphi'_k(s_j), \sum_{k=1}^n\texttt{c}^2_k\varphi'_k(s_j)\right)\right\|_{\R^2}^2-1=0,\\[2mm]
                     & \int_{0}^{1} \sum_{i=1}^n\texttt{a}_i\varphi_i(r)dr=0,\\[2mm]
                     & \frac{1}{\mu_t^2} \left(\sum_{i=1}^n\texttt{a}_i\varphi_i(s_j)\right)^2 + \frac{1}{\mu_n^2 p^2}\texttt{v}(s_j)^4\left\| \left(\sum_{k=1}^n\texttt{c}^1_k\varphi_k''(s_j), \sum_{k=1}^n\texttt{c}^2_k\varphi_k''(s_j)\right) \right\|_{\R^2}^2-g^2\leq 0,\,\\[2mm]
                     &  \|p\left(\sum_{k=1}^n\texttt{c}^1_k\varphi_k(s_j), \sum_{k=1}^n\texttt{c}^2_k\varphi_k(s_j)\right)-\alpha(s_j)\|_{\R^2}^2-\kappa(s_j)^2 \leq 0,\, \\[2mm]
                     &  \|p\left(\sum_{k=1}^n\texttt{c}^1_k\varphi_k(s_j), \sum_{k=1}^n\texttt{c}^2_k\varphi_k(s_j)\right)-\beta(s_j)\|_{\R^2}^2-\kappa(s_j)^2 \leq 0, \,\\[2mm]
                     & p\geq 0,\,\\[2mm]
                     & v_0\geq 0,

  \end{array}
\end{equation}
donde las restricciones (salvo la segunda) deben cumplirse para cada $j=1,\ldots l$.

Con el problema final en dimensión finita \eqref{pdisfinal} planteado, podemos demostrar que dicho problema tiene solución.

\subsection{Existencia de solución para el problema discretizado}\label{existsoldiscretizado}

Para probar la existencia de solución de el problema \eqref{pdisfinal}, observemos que la función $\texttt{f}$ claramente es continua, por lo tanto si para un punto factible $\texttt{x}_0$ dado, el conjunto $\{\texttt{x}\in S\mid \texttt{f}(\texttt{x})\leq \texttt{f}(\texttt{x}_0)\}$ es compacto, entonces el Corolario \ref{solglobalfinita} nos garantiza la existencia de solución. Veamos entonces que este conjunto es compacto. Para simplificar notación, sea 
$$R=\{\texttt{x}\in\R^{3n+2}\mid \texttt{f}(\texttt{x})\leq \texttt{f}_0\},$$
donde $\texttt{f}_0=\texttt{f}(\texttt{x}_0)$.
\begin{teo}
El conjunto $S\cap R$ es compacto no vacío.
\end{teo}
\begin{dem}
Por la continuidad de las funciones y el hecho que $\texttt{x}_0$ es factible, los conjuntos $S$ y $R$ son cerrados y no vacíos. Ya que $S\cap R$ es cerrado y contiene a $\texttt{x}_0$, basta ver que es acotado.

Supongamos que existe un $\texttt{x}^k=((\texttt{c}^{1})^k,(\texttt{c}^{2})^k,\texttt{a}^k,p_k,(v_0)_k)\in\R^{3n+2}$ que pertenece al conjunto $S\cap R$  tal que $\|\texttt{x}^k\|_{\R^2}\to\infty$, y llamemos $\eta_k:=\|\texttt{x}^k\|_{\R^2}$. Luego, tomando subsucesión si fuera necesario, suponemos que $\frac{\texttt{x}^k}{\eta_k}\to\texttt{x}_s=(\texttt{c}^1_s,\texttt{c}^2_s,\texttt{a}_s,p_s,(v_0)_s)$ con $\|\texttt{x}_s\|_{\R^2}=1$. Veamos que esto genera un absurdo.

Primero, por la restricción de la elipse de fricción, tenemos que para todo $j$, $$\frac{1}{\mu_t^2} \left(\sum_{i=1}^n\texttt{a}^k_i\varphi_i(s_j)\right)^2 + \frac{1}{\mu_n^2 p_k^2}\texttt{v}^k(s_j)^4\left\| \left(\sum_{i=1}^n(\texttt{c}^{1})^k_i\varphi_i''(s_j), \sum_{i=1}^n(\texttt{c}^{2})^k_i\varphi_i''(s_j)\right) \right\|_{\R^2}^2-g^2\leq 0,$$
entonces, como el segundo miembro es mayor o igual que cero, vale que para todo $j$,
\begin{equation}\label{restriccionfriccionmodificada}
    \left(\sum_{i=1}^n\texttt{a}^k_i\varphi_i(s_j)\right)^2\leq g^2\mu_t^2,
\end{equation}



luego si divido ambos miembros por $\eta_k^2$ y tomo límite, se obtiene
$$\sum_{i=1}^n(\texttt{a}_s)_i\varphi_i(s_j)=0,$$
finalmente por la independencia lineal de las funciones base, podemos asegurar que
\begin{equation}\label{aigualacero}
    (\texttt{a}_s)_i=0 \ \forall i=1,\ldots,n.
\end{equation}

Notemos que, para todo $s$,
\begin{equation}\label{vtiendeav0}
   \frac{\texttt{v}^k(s)}{\eta_k}=\sqrt{2\frac{p_k}{\eta_k}\int_0^s \sum_{i=1}^n\frac{\texttt{a}^k_i}{\eta_k}\varphi_i(r)dr +\left(\frac{(v_0)_k}{\eta_k}\right)^2} \mathop{\longrightarrow}_{k\to\infty}(v_0)_s, 
\end{equation}
donde se usó \eqref{aigualacero}.




Ahora, observemos que como se debe cumplir que $\texttt{f}(\texttt{x}^k)\leq \texttt{f}_0$, esto es, $\int_0^1\frac{p_k}{\texttt{v}^k(s)}ds\leq\texttt{f}_0$, si multiplico y divido ambos miembros por $\eta_k$ y uso \eqref{vtiendeav0}, obtengo $$\int_0^1\frac{\frac{p_k}{\eta_k}}{\frac{\texttt{v}^k(s)}{\eta_k}}ds\mathop{\longrightarrow}_{k\to\infty}\int_0^1\frac{p_s}{(v_0)_s} ds =\frac{p_s}{(v_0)_s},$$
y por lo tanto,
\begin{equation}\label{desigualdadf0v0}
    \texttt{f}_0\geq\frac{p_s}{(v_0)_s}.
\end{equation}

Luego utilizaremos el resultado anterior. Consideremos ahora la ecuación \eqref{restriccionfriccionmodificada}, entonces, de una manera análoga a la utilizada para obtener esta fórmula, se tiene que para todo $j$,
\begin{equation*}
    \left\| \left(\sum_{i=1}^n(\texttt{c}^{1})^k_i\varphi_i''(s_j), \sum_{i=1}^n(\texttt{c}^{2})^k_i\varphi_i''(s_j)\right) \right\|_{\R^2}^2\leq g^2\mu_n^2\frac{p_k^2}{\texttt{v}^k(s_j)^4},
\end{equation*}
y podemos aplicar raíz cuadrada a ambos lados de la desigualdad para obtener
\begin{equation}\label{desigualdadd2c}
    \left\| \left(\sum_{i=1}^n(\texttt{c}^{1})^k_i\varphi_i''(s_j), \sum_{i=1}^n(\texttt{c}^{2})^k_i\varphi_i''(s_j)\right) \right\|_{\R^2}\leq g\mu_n\frac{p_k}{\texttt{v}^k(s_j)^2}.
\end{equation}

Ahora, consideremos la función $\hat{c}^k:(-\infty,\infty)\arr\R^2$ tal que
\[
 \hat{c}^k(s):=\left(\sum_{i=1}^n(\texttt{c}^{1})^k_i\varphi_i(s), \sum_{i=1}^n(\texttt{c}^{2})^k_i\varphi_i(s)\right).
\]
Note que esta función es infinitamente diferenciable y periódica en $[0,1]$. Como la función $s\mapsto\|\hat{c}^k(s)\|_{\R^2}$ es continua y $[0,1]$ compacto, entonces esta funciçon alcanza su máximo en $s_0\in[0,1]$. Por lo tanto $s_0$ es maximizador en toda la recta real. Entonces podemos aplicar la Proposición \ref{cotainfcurvatura}, y usando que la distancia de la curva al origen es siempre menor o igual que la de $\beta$ al origen, obtenemos que 
$$\left\| \left(\sum_{i=1}^n(\texttt{c}^{1})^k_i\varphi_i''(s_0), \sum_{i=1}^n(\texttt{c}^{2})^k_i\varphi_i''(s_0)\right) \right\|_{\R^2}\geq\frac{1}{\left\| \hat{c}^k(s_0) \right\|_{\R^2}}\geq\frac{1}{\|\beta(s_0)\|_{\R^2}}.$$
Luego como la curvatura es continua, dado un $\varepsilon>0$ existe un entorno abierto $U$ de $s_0$ tal que para todo $s\in U$
$$\left\| \left(\sum_{i=1}^n(\texttt{c}^{1})^k_i\varphi_i''(s), \sum_{i=1}^n(\texttt{c}^{2})^k_i\varphi_i''(s)\right) \right\|_{\R^2}>\frac{1}{\|\beta(s_0)\|_{\R^2}+\varepsilon}.$$
Entonces tomo un $n$ suficientemente grande, de manera que exista un $j_0$ tal que $s_{j_0}\in U$. Por lo tanto, vale
\begin{equation}\label{desigcurvatura}
    \left\| \left(\sum_{i=1}^n(\texttt{c}^{1})^k_i\varphi_i''(s_{j_0}), \sum_{i=1}^n(\texttt{c}^{2})^k_i\varphi_i''(s_{j_0})\right) \right\|_{\R^2}\geq\frac{1}{\|\beta(s_0)\|_{\R^2}+ \varepsilon}.
\end{equation}






Ahora, si multiplicamos ambos lados de la desigualdad \eqref{desigcurvatura} por $\eta_k$, utilizando \eqref{desigualdadd2c} y tomando límite se tiene que
\begin{equation}\label{desetapv0}
    \infty\mathop{\longleftarrow}_{k\to\infty}\frac{\eta_k}{\|\beta(s_0)\|_{\R^2}+\varepsilon }\leq g\mu_n\frac{\frac{p_k}{\eta_k}}{\left(\frac{\texttt{v}^k(s_{j_0})}{\eta_k}\right)^2}\mathop{\longrightarrow}_{k\to\infty}g\mu_n\displaystyle\frac{p_s}{(v_0)^2_s},
\end{equation}
donde \eqref{vtiendeav0} es usada.

Supongamos que $p_s>0$, entonces por \eqref{desetapv0}, $(v_0)_s=0$. Pero si miramos la ecuación \eqref{desigualdadf0v0} esto es absurdo. Por lo tanto $p_s=0$. Ahora si $(v_0)_s>0$ obtenemos otro absurdo debido a \eqref{desetapv0}, entonces $(v_0)_s$ debe ser $0$. El resultado fundamental para terminar la demostración es el \texttt{Teorema de la desigualdad isoperimétrica} (ver Teorema \ref{desigualdadisoperimetrica} del Apéndice), por este teorema el cuadrado la longitud $p_k$ de la trayectoria está acotado inferiormente por $4\pi$ veces el área de la región limitada por esta trayectoria. Por otro lado, como el área que encierra la curva interior $\alpha$ ($B_\alpha$) siempre es menor o igual que el de la trayectoria, obtenemos que
\begin{equation*}
    p_k^2\geq 4\pi B_\alpha>0.
\end{equation*}
Y esto nos dice que $\displaystyle \eta_kp_k\mathop{\longrightarrow}_{k\to\infty}\infty$, luego, considerando la restricción de permanencia en pista (que involucra a $\alpha$), tenemos (para todo $j$)
\begin{equation*}
    \left\| \left(\sum_{i=1}^n\frac{(\texttt{c}^{1})^k_i}{\eta_k}\varphi_i(s_j), \sum_{i=1}^n\frac{(\texttt{c}^{2})^k_i}{\eta_k}\varphi_i(s_j)\right)-\frac{\alpha(s_j)}{\eta_kp_k} \right\|_{\R^2} \leq \frac{\kappa(s_j)}{\eta_kp_k},
\end{equation*}
entonces tomando límite obtenemos que
\[
 \left\|\left(\sum_{i=1}^n(\texttt{c}^1_s)_i\varphi_i(s_j), \sum_{i=1}^n(\texttt{c}^2_s)_i\varphi_i(s_j)\right)\right\|_{\R^2}=0.
\]

% \begin{equation*}
%     0\mathop{\longleftarrow}_{k\to\infty}\frac{\kappa(s_j)}{\eta_kp_k}\geq \left\| \left(\sum_{i=1}^n\frac{(\texttt{c}^{1})^k_i}{\eta_k}\varphi_i(s_j), \sum_{i=1}^n\frac{(\texttt{c}^{2})^k_i}{\eta_k}\varphi_i(s_j)\right)-\frac{\alpha(s_j)}{\eta_kp_k} \right\|_{\R^2}\mathop{\longrightarrow}_{k\to\infty}\left\|\left(\sum_{i=1}^n(\texttt{c}^1_s)_i\varphi_i(s_j), \sum_{i=1}^n(\texttt{c}^2_s)_i\varphi_i(s_j)\right)\right\|_{\R^2}.
% \end{equation*}

Por lo tanto, debido a que las funciones base son linealmente independientes, resulta $c^1_s=0$ y $c^2_s=0$. Como sabemos que $a_s=0$, $p_s=0$ y $(v_0)_s=0$, vale que $x_s=0$ y esto es un absurdo pues la norma de $x_s$ debía ser igual a $1$.
\end{dem}

Luego, para utilizar un método numérico y encontrar la solución del problema, necesitaremos calcular las derivadas parciales ya en dimensión finita, de allí, la siguiente sección.

\subsection{Derivadas}

\textbf{ESCRIBIR LAS DERIVADAS FINALES}

\begin{eqnarray*}
\frac{\partial \texttt{f}}{\partial \texttt{c}^1_k}(\texttt{x})&=&\langle f'(\hat x) , (\varphi_k,0,0,0,0)\rangle\\
\frac{\partial \texttt{f}}{\partial \texttt{c}^2_k}(\texttt{x})&=&\langle f'(\hat x) , (0,\varphi_k,0,0,0)\rangle\\
\frac{\partial \texttt{f}}{\partial \texttt{a}_k}(\texttt{x})&=&\langle f'(\hat x) , (0,0,\varphi_k,0,0)\rangle\\
\frac{\partial \texttt{f}}{\partial p}(\texttt{x})&=&\langle f'(\hat x) , (0,0,0,1,0)\rangle\\
\frac{\partial \texttt{f}}{\partial v_0}(\texttt{x})&=&\langle f'(\hat x) , (0,0,0,0,1)\rangle\\
(\mbox{Entonces queda, por ejemplo} \quad\frac{\partial \texttt{f}}{\partial \texttt{a}_k}(\texttt{x})&=&\int^1_0-\frac{p^2}{\texttt{v}(s)^3}\int^s_0\varphi_k(r)dr ds)\\
(\texttt{G}_l)_i(\texttt{x})&=&\langle G_l(x) , \varphi_i\rangle\\
\frac{\partial(\texttt{G}_l)_i}{\partial\texttt{c}^2_k}(\texttt{x})&=&\langle G'_l(x)(0,\varphi_k,0,0,0) , \delta_{s_i}\rangle\\
(\mbox{Y queda, por ejemplo,}\frac{\partial(\texttt{G}_7)_i}{\partial p}(\texttt{x})&=&\frac{4\texttt{v}(i)}{p^2\mu_n^2}\left\|\frac{d^2\texttt{c}(i)}{ds^2}\right\|_{\R^2}^2\sum_{j=0}^n\texttt{a}_j\int^1_0\varphi_j(r)dr
\end{eqnarray*}

\begin{eqnarray*}
G_7'(x)(\tilde{x})&=&\left.\frac{\partial}{\partial\alpha}\right|_0 G(x+\alpha x)\\
=\frac{(a+\alpha\tilde{a})^2}{\mu_t^2}+\frac{1}{(p+\alpha\tilde{p})^2\mu_n^2}F(\alpha)\left\|\frac{d^2c}{ds^2}+\alpha\frac{d^2\tilde{c}}{ds^2}\right\|_{\R^2}^2&-&\left.g^2\right|_0\\
(\mbox{donde }F(\alpha)=(2(p+\alpha\tilde{p})\int_0^s a+\alpha\tilde{a}dr+(v_0&+&\alpha\tilde{v}_0)^2)^2\mbox{ y}\\
F'(\alpha)=2(2(p+\alpha\tilde{p})\int_0^s) a+\alpha\tilde{a}dr+(v_0+\alpha\tilde{v}_0)^2)[2\tilde{p}\int_0^s a+\alpha\tilde{a}dr&+&2\int_0^s\tilde{a}dr(p+\alpha\tilde{p})+2(v_0+\alpha\tilde{v}_0)\tilde{v}_0])\\
=\frac{2(a+\alpha\tilde{a})\tilde{a}}{\mu_t^2}+\frac{1}{\mu_n^2}[\frac{-2(p+\alpha\tilde{p})\tilde{p}}{(p+\alpha\tilde{p})^4}&F(\alpha)&\left\|\frac{d^2c}{ds^2}+\alpha\frac{d^2\tilde{c}}{ds^2}\right\|_{\R^2}^2\\
+\frac{1}{(p+\alpha\tilde{p})^2}(F'(\alpha)\left\|\frac{d^2c}{ds^2}+\alpha\frac{d^2\tilde{c}}{ds^2}\right\|_{\R^2}^2&+&\left.F(\alpha)2\left\la \frac{d^2c}{ds^2}+\alpha\frac{d^2\tilde{c}}{ds^2} , \frac{d^2\tilde{c}}{ds^2}\right\ra)]\right|_0\\
=\frac{2a\tilde{a}}{\mu_t^2}+\frac{1}{\mu_n^2}\left[-\frac{2p\tilde{p}}{p^4}(2p\int_0^s adr\right.&+&v_0^2)^2\left\|\frac{d^2c}{ds^2}\right\|_{\R^2}^2\\
+\frac{1}{p^2}\left((2(2p\int_0^s adr+v_0^2))(2\tilde{p}\int_0^s adr+2p\int_0^s\tilde{a}dr+2v_0\tilde{v}_0)\left\|\frac{d^2c}{ds^2}\right\|_{\R^2}^2\right.&+&\left.\left.(2p\int_0^s adr+v_0^2)^22\left\la \frac{d^2c}{ds^2} , \frac{d^2\tilde{c}}{ds^2}\right\ra\right)\right]\\
\end{eqnarray*}

\begin{itemize}
 \item$(\tilde{c})\arr\frac{v^4}{\mu_n^2p^2}2\left\la \frac{d^2c}{ds^2} , \frac{d^2\tilde{c}}{ds^2}\right\ra$
 \item$(\tilde{a})\arr\frac{2a\tilde{a}}{\mu_t^2}+\frac{4v^2}{\mu_n^2p^2}p\int_0^s\tilde{a}\left\|\frac{d^2c}{ds^2}\right\|_{\R^2}^2$
 \item$(\tilde{p})\arr\frac{1}{\mu_n^2}\left[-\frac{2p\tilde{p}}{p^4}v^4\left\|\frac{d^2c}{ds^2}\right\|_{\R^2}^2+\frac{4v^4}{p^2}\tilde{p}\int_0^s adr\left\|\frac{d^2c}{ds^2}\right\|_{\R^2}^2\right]$
 \item$(\tilde{v}_0)\arr\frac{4v^2v_0}{\mu_n^2p^2}\left\|\frac{d^2c}{ds^2}\right\|_{\R^2}^2$
\end{itemize}
    
Las demás derivadas parciales de la $(\textbf{G}_l)_i$ se escriben análogamente moviendo de lugar el $\varphi_k$



\chapter{Resultados computacionales}

Con el objetivo de iniciar con una pista simple, utilizamos el algoritmo SQP para encontrar la trayectoria óptima en un circuito circular, es decir, tanto la calzada interna como la externa son ambas circunferencias de, en este caso radios $1$ y $1.5$ respectivamente \textbf{escribir demas condciones inicales}. Así obtenemos

\textbf{graficas pista cicular}

Se puede visualizar que la trayectoria óptima en este caso, es recorrer la calzada interna. Tal vez este resultado pueda parecer un poco trivial, procedemos entonces a variar un poco más los circuitos.

Por ejemplo, si ahora la calzada interna y externa son rectangulares, se logra lo siguiente,

\textbf{grafica pista cuadrada 1}

Se puede observar ahora, que seguir la calzada interna, ya no es la trayectoria óptima, sino, lo optimal es una trayectoria suave (semejante a una circunferencia) que sólo interseca la calzada interna en los vértices de la misma.

\textbf{la otra pista cuadrada}

Finalmente, utilizamos el algoritmo para encontrar la trayectoria óptima de un circuito más complejo, el de una pista de carreras. El resultado obtenido fue

\textbf{grafica pista de carrera}

Observando la gráfica de la trayectoria, podemos ver que la misma, intenta mantener cierta suavidad constantemente, de manera que pueda lograr la mayor aceleración posible, teniendo en cuenta la fricción de la pista.

\chapter{Conclusión}

De acuerdo a lo realizado en este trabajo, somos (plural está bien?) capaces de concluir que se logró el objetivo de plantear y resolver el problema de optimización propuesto, es decir, encontramos una manera de obtener una trayctoria óptima, en un circuito cerrado. Además el código programado también nos provee la aceleración tangencial y la rapidez. También en este trabajo se realizo una introducción a la teoría de existencia de soluciones para problemas de optimización en dimensión infinita y existencia de multiplicadores de Lagrange tanto en dimensión infinita como en dimensión finita. Además de una descripción del método computacional SQP, utilizado para resolver el problema de optimización.

El trabajo puede ser utilizado para cualquier circuito de carreras, sólo es necesario dar las coordenadas de las pistas como vimos en \textbf{referenciar}. Sin embargo puede tener diversos usos, como por ejemplo optimizar el recorrido en líneas de producción de fábricas por ejemplo, sólo habría que modificar el problema planteando la correspondiente ecuación de restricción de fricción (ya no necesariamente sería una elipse de fricción).

El problema de encontrar una trayectoria óptima es muy estudiado, en particular hay trabajos que se concentran en circuitos automovilísticos, tomando en cuenta muchas más restricciones físicas, donde el vehículo ya no es puntual, sino que es un automóvil. Artículos del estilo pueden verse por ejemplo en \cite{xiong2010racing} y \cite{van2015optimal}. 







\chapter{Apéndice}



\section{Códigos para MATLAB}


%\begin{defi}
% La funcion Lagrangiana aumentada asociada a ($\tilde{P}$) es
% 
% $L_c:\R^n\times\R^m\times\R^r\arr\R$ tal que 
% $$L_c(x,\lambda,\mu)=f(x)+\la\lambda , h(x)\ra %+\frac{c}{2}\|h(x)\|^2+\frac{1}{2c}\sum_{j=1}^r(\mbox{máx}\{0,\mu_j+cg_j(x)\}^2-\mu_j^2)$%$.
%\end{defi}

%\textbf{Método del Lagrangiano aumentado}

%Dado $\lambda^0, \mu^0, c_0^0>0$; genero $\{x^k\}, \{\lambda^k\}, \{\mu^k\},\{c_k\}$

%PASO 1. Tomo $x^{k+1}$ solución de 
%\[
%\left\{\begin{array}{rl}
%\displaystyle \mathop{\mathrm{minimizar}}_{x\in\Omega} & \displaystyle %L_{c_k}(x,\lambda^k, \mu^k)
%\end{array}\right.
%\].
%PASO 2. Defino $$\lambda^{k+1}= \lambda^k + c_k h(x^{k+1})$$ y
%$$\mu^{k+1}=\mbox{máx}\{0,\mu^k+c_k g(x^{k+1})\}$$.

%ASO 3. Tomo $c_{k+1}\geq c_k$.

%\begin{prop}
 %Asuma f,g y h continuas. $\Omega\subset\R^n$ cerrado y $X$ no vacío. Consideremos las %sucesiones $\{x^k\}, \{\lambda^k\},\{\mu^k\} y \{c_k\}$ generadas por el método del %Lagrangiano aumentado tal que $x^{k+1}$ es solución global, $\{\lambda^k\},\{\mu^k\}$ %acotadas y $c_k\arr\infty$ Entonces todo punto límite de $\{x^k\}$ es un mínimo global %de $\tilde{P}$.
%\end{prop}


















%\nocite{*}
\bibliographystyle{abbrv}%{plain}
\bibliography{bibliografia}  

\end{document}